\documentclass[12pt, oneside]{article}

\usepackage{xparse}

\usepackage{geometry}
\geometry{letterpaper}
\usepackage[parfill]{parskip}
\usepackage{float}

% Math packages
\usepackage{amssymb}
\usepackage{amsmath}
\usepackage{amsthm}
\usepackage{mathalpha}

% Tables
\usepackage{array}
\usepackage{siunitx}
\usepackage{booktabs}
\usepackage{afterpage}

% Hyperlinks
\usepackage{hyperref}

% References
\numberwithin{equation}{section}
\usepackage[numbers]{natbib}

% Theorem environments
\theoremstyle{definition}
\newtheorem{definition}{Definition}[section]
\newtheorem{proposition}{Proposition}[section]
\newtheorem{axiom}{Axiom}[section]

\theoremstyle{plain}
\newtheorem{theorem}{Theorem}[section]
\newtheorem{lemma}[theorem]{Lemma}
\newtheorem{corollary}[theorem]{Corollary}
\newtheorem{conjecture}[theorem]{Conjecture}
\newtheorem{example}[theorem]{Example}

\theoremstyle{remark}
\newtheorem*{remark}{Remark}

\title{A Symbolic--Affine Framework for Density-One Collatz Descent}
\author{Wayne Brassem}

\begin{document}

% Document commands which provide hyperlinks to OEIS sequences referenced herein
\NewDocumentCommand{\OEIS}{m}{\href{https://oeis.org/#1}{#1}}

\maketitle

\begin{abstract}

We study finite stopping time for the Collatz map through the fully
accelerated transformation
\[
f(x)=\frac{3x+1}{2^{v_2(3x+1)}},
\]
acting on the odd positive integers.

Finite accelerated trajectories are encoded by \emph{division words},
which record the successive \(2\)-adic valuations encountered along an
orbit. Each division word determines a formal affine map
\[
F_\omega(x)
=
\frac{3^{|\omega|}x+c(\omega)}{2^{D(\omega)}}.
\]
For every nonempty word with positive entries, exact realization is
equivalent to a single affine congruence modulo
\(2^{D(\omega)+1}\). Consequently, its realizations form a unique dyadic
cylinder of natural density \(2^{-(D(\omega)+1)}\), and symbolic prefix
extension corresponds exactly to cylinder containment.

We organize these cylinders into a residual forest according to the first
prefix at which the affine slope becomes strictly less than \(1\). Residual
words have affine slope strictly greater than \(1\) at every nonempty prefix,
whereas first-contraction words acquire subunit affine slope for the first
time at their terminal position. The resulting residual sets satisfy an
exact finite-stage decomposition into first-contraction layers and continued
residual survival.

To prove that the residual mass vanishes, we introduce the auxiliary weight
\[
\left(\frac38\right)^{D(\omega)},
\]
whose total residual weight contracts by a factor at most \(3/5\) at each
symbolic step. Comparison with the true dyadic cylinder weights, together
with
\[
B(m+5)\le B(m)+8,
\]
yields a five-step envelope contraction with factor
\[
\left(\frac43\right)^8\left(\frac35\right)^5
=
\frac{65536}{84375}
<1.
\]
Hence the natural densities of the finite-depth residual sets tend to zero.

A uniform cutoff on each first-contraction layer converts affine contraction
into genuine descent apart from finitely many exceptional values.

Since the initial residual class is exactly \(3\bmod4\), while positive even
integers descend after one ordinary Collatz step and odd integers congruent
to \(1\bmod4\), apart from \(1\), descend after one fully accelerated step,
it follows that the set of positive integers having an ordinary Collatz
iterate below their starting value has natural density one.

The proof is deterministic, uses finite-stage natural-density accounting,
and is independently formalized in Lean~4.

\end{abstract}

\newpage
\tableofcontents

\newpage
% ============================================================
\section{Introduction}
% ============================================================

The \(3x+1\) problem, commonly known as the Collatz conjecture, concerns
iteration of the map
\[
T(x)=
\begin{cases}
\dfrac{x}{2}, & x \text{ even},\\[6pt]
3x+1, & x \text{ odd},
\end{cases}
\]
on the positive integers. The conjecture asserts that every positive integer
eventually reaches \(1\). Despite its elementary formulation, the problem
has resisted proof for decades; historical accounts and surveys of the
problem and its many equivalent formulations are given by Lagarias
\cite{Lagarias2010}. Extensive computation has also verified convergence
over large finite ranges \cite{OliveiraESilva1999}.

A classical weakening of the conjecture asks only whether an orbit eventually
falls below its starting value. The least positive iterate at which this
occurs is commonly called the \emph{stopping time}. Terras
\cite{Terras1976} and, independently, Everett \cite{Everett1977} established
that finite stopping time occurs for almost every positive integer in the
sense of natural density.

Several ideas underlying such finite-stopping arguments are already present
in the classical work. Terras encodes finite trajectories by parity data and
expresses the resulting iterates in affine form, with a multiplicative
coefficient determined by that finite symbolic data together with an
additive remainder. The associated \emph{coefficient stopping time} records
the first depth at which the multiplicative coefficient becomes less than
\(1\). The finite parity data also determine arithmetic progressions of
starting values, providing a residue-class structure from which density
statements can be obtained; see \cite{Terras1976,Lagarias2010}. Thus affine
trajectory representations, finite symbolic coding, residue-class
organization, and the density-one finite-stopping theorem itself are
classical ingredients.

Subsequent work has obtained substantially stronger control of orbit minima
in other density settings. Tao \cite{Tao2019} proved that, in logarithmic
density, almost every Collatz orbit attains almost bounded values: for every
function \(g(N)\to\infty\), the minimum value attained by the orbit of \(N\)
is at most \(g(N)\) for logarithmic-density almost every \(N\).

The purpose of the present work is therefore not to introduce affine
trajectory coding or to claim a new density-one finite-stopping theorem.
Rather, it develops a self-contained deterministic formulation of the
finite-stopping argument in fully accelerated coordinates and organizes the
proof around exact realization classes and a residual-cylinder
decomposition. The odd part of the dynamics is encoded using the fully
accelerated map
\[
f(x)
=
\frac{3x+1}{2^{v_2(3x+1)}},
\]
which maps odd positive integers to odd positive integers by removing all
factors of \(2\) following each \(3x+1\) step.

A finite trajectory under \(f\) determines a \emph{division word}
\[
\omega=(d_0,\dots,d_{m-1}),
\qquad
d_j=v_2\!\left(3f^j(x)+1\right),
\]
recording the exact powers of \(2\) removed at successive accelerated
steps. Write
\[
D(\omega)=d_0+\cdots+d_{m-1}
\]
for its total division count. Symbolic composition associates to every
division word the formal affine map
\[
F_\omega(x)
=
\frac{3^{|\omega|}x+c(\omega)}
     {2^{D(\omega)}}.
\]
This is the fully accelerated analogue of the classical affine description
of finite Collatz trajectories. It is algebraic and is defined independently
of whether the prescribed word is realized by an actual accelerated
trajectory.

In these coordinates the arithmetic realization problem has an exact
congruential form. For an admissible division word---that is, a nonempty
word whose entries are positive---a natural number \(x\) realizes
\(\omega\) if and only if
\[
3^{|\omega|}x+c(\omega)
\equiv
2^{D(\omega)}
\pmod{2^{D(\omega)+1}}.
\]
Consequently, the realization set of \(\omega\) is a unique residue class
modulo
\[
2^{D(\omega)+1}.
\]
In particular, every admissible finite division word is realizable. The
additional factor of \(2\) in the modulus distinguishes exact accelerated
realization from the weaker condition that the formal affine expression is
merely integral.

These realization classes form nested dyadic cylinders. If \(C_\omega\)
denotes the realization cylinder of \(\omega\), then
\[
\operatorname{dens}(C_\omega)
=
2^{-(D(\omega)+1)},
\]
and symbolic prefix extension corresponds exactly to cylinder containment.
Prefix-incomparable words determine disjoint cylinders. Thus the classical
residue-class organization of finite trajectories takes, in the present
accelerated coordinates, the form of a nested dyadic cylinder geometry.

The affine slope of a division word is
\[
\lambda_\omega
=
\frac{3^{|\omega|}}{2^{D(\omega)}}.
\]
Let
\[
B(m)
=
\lfloor m\log_2 3\rfloor+1
\]
denote the least integer \(k\) for which
\[
3^m<2^k.
\]
Then
\[
\lambda_\omega<1
\quad\Longleftrightarrow\quad
D(\omega)\ge B(|\omega|).
\]
This is the accelerated-coordinate counterpart of the classical
coefficient-stopping comparison: it concerns the multiplicative affine
coefficient rather than pointwise descent itself.

A subunit affine slope does not by itself imply descent for every
realization, because the positive affine constant may dominate for small
initial values. Nevertheless,
\[
F_\omega(x)<x
\]
for every sufficiently large realization of a fixed subunit-slope word.
Hence only finitely many points of such a realization cylinder can fail to
lie below their starting value at the endpoint of the encoded trajectory
segment, and removing those points does not change the cylinder's natural
density.

The global argument is organized by the first prefix at which this affine
threshold is reached. An admissible word is called \emph{residual} when every
nonempty prefix remains above unit affine slope, and a
\emph{first-contraction word} when its proper prefixes are residual but its
terminal prefix has subunit slope. Although the full family of admissible
words is infinite at every positive symbolic length, the residual family at
each fixed depth is finite.

If \(\mathcal R_m\) denotes the union of the realization cylinders of
residual words of length \(m\), and \(\mathcal F_{m+1}\) the corresponding
first-contraction exit layer, then
\[
\mathcal R_m
=
\mathcal F_{m+1}
\sqcup
\mathcal R_{m+1}.
\]
This exact finite-stage decomposition provides the conservation law used
throughout the residual-forest argument.

Writing
\[
M_m
=
\operatorname{dens}(\mathcal R_m)
\]
for the residual mass, the quantitative task is to prove
\[
M_m\longrightarrow0.
\]
The proof introduces the auxiliary weight
\[
\left(\frac38\right)^{D(\omega)},
\]
whose total weight satisfies the uniform one-step bound
\[
W_{m+1}\le\frac35\,W_m.
\]
Comparison with the true dyadic cylinder weights, together with
the elementary threshold estimate
\[
B(m+5)\le B(m)+8,
\]
yields a five-step envelope contraction with factor
\[
\left(\frac43\right)^8
\left(\frac35\right)^5
=
\frac{65536}{84375}
<1.
\]
It follows that the unresolved residual mass tends to zero.

The final step converts affine first contraction into genuine dynamical
descent. Although a fixed first-contraction layer may contain infinitely
many terminal words, its residual parent level is finite. The affine
concatenation formula therefore supplies a uniform finite cutoff above
which every point in the entire exit layer descends. At any fixed residual
depth, the set of non-descending points is consequently contained in the
remaining residual set together with a finite exceptional set. Taking the
residual depth first and then the interval size shows that initially
residual non-descent has natural density zero.

The initial residual set itself has the simple description
\[
\mathcal R_1
=
\{x:x\equiv3\pmod4\}.
\]
Positive even integers descend immediately under \(T\), while odd integers
congruent to \(1\pmod4\), apart from \(1\), descend after one fully
accelerated step. Accelerated iterates occur along the ordinary Collatz
orbit, so accelerated descent transfers to ordinary descent. Thus the
residual argument handles the remaining nontrivial congruence class and
recovers the classical density-one finite-stopping theorem:
\[
\operatorname{dens}
\left\{
x\in\mathbb N_{>0}:
\exists n,\;T^n(x)<x
\right\}
=
1.
\]

The resulting proof is organized through the sequence
\[
\begin{aligned}
\text{division words}
&\longrightarrow
\text{affine realization}
\longrightarrow
\text{dyadic cylinders}
\\
&\longrightarrow
\text{residual forest}
\longrightarrow
\text{residual mass collapse}
\longrightarrow
\text{eventual descent}.
\end{aligned}
\]
This organization should be viewed as a fully accelerated, exact-arithmetic
reformulation of classical finite-stopping machinery rather than as a claim
of priority for affine or residue coding. Its distinguishing role in the
present proof is to isolate the unresolved prefixes into finite residual
levels and control their total dyadic mass through the residual-forest
decomposition. No probabilistic model of trajectory independence is used,
and the density argument proceeds through finite-stage decompositions and
ordinary natural-density limits.

The result concerns descent below the starting value, not convergence to
\(1\). It does not assert that every positive starting value has finite
stopping time, nor does it prove that every Collatz trajectory converges.
The density-zero complement of the descent set may still be nonempty or
infinite.

\paragraph{Formal verification and reproducibility.}
The Lean~4 formalization accompanying this manuscript is available at
\begin{center}
\url{https://github.com/wbrassem/collatz-density-one-descent}
\end{center}
The culminating formal theorem is
\begin{center}
\nolinkurl{Collatz.ordinary_collatz_descent_has_natural_density_one}
\end{center}
The repository contains the Lean project, verification harness,
continuous-integration workflow, and manuscript source.

\paragraph{Structure of the Paper.}
The argument proceeds as follows:
\begin{itemize}
\item Section~\ref{sec:division-counts} introduces division words,
realization, and total division counts.
\item Section~\ref{sec:affine-representation} develops the formal affine
representation and the dyadic--ternary slope threshold.
\item Section~\ref{sec:exact-realizability-residue} proves the exact
realization criterion and the unique residue class associated with each
admissible word.
\item Section~\ref{sec:cylinder-geometry} develops generic dyadic-cylinder
geometry, counting, natural density, and the realization-cylinder forest.
\item Section~\ref{sec:affine-contraction-descent} relates subunit affine
slope to pointwise descent and proves density preservation after removing
the finite affine exceptional set.
\item Section~\ref{sec:residual-forest} introduces residual and
first-contraction words and establishes the finite-stage forest and mass
decompositions.
\item Section~\ref{sec:residual-mass-collapse} proves that the residual mass
tends to zero by means of the auxiliary weighted mass and the five-step
envelope contraction.
\item Section~\ref{sec:density-one-descent} converts first contraction into
genuine descent, identifies the initial residual congruence class, and
transfers the accelerated result to the ordinary Collatz map.
\end{itemize}


% ============================================================
\section{Conventions}
\label{sec:conventions}
% ============================================================

Throughout,
\[
\mathbb N=\{0,1,2,\ldots\},
\qquad
\mathbb N_{>0}=\{1,2,3,\ldots\}.
\]

The classical Collatz conjecture concerns iteration on
\(\mathbb N_{>0}\). The ordinary Collatz map is therefore regarded as a
dynamical map on positive integers, while the fully accelerated Collatz map
acts dynamically on the odd positive integers.

For the arithmetic and symbolic constructions used in this paper, however,
it is convenient to work in the ambient set \(\mathbb N\). Generic residue
classes, dyadic cylinders, symbolic data, exponents, indices, and counting
functions are therefore allowed to involve \(0\).

Natural density is computed using the initial intervals
\[
[0,N)=\{0,1,\ldots,N-1\}.
\]
Thus a subset of \(\mathbb N_{>0}\) is also regarded as a subset of
\(\mathbb N\) for purposes of density; the omission of the single point
\(0\) has no effect on its natural density.

Division words may contain the value \(0\). Positivity conditions required
for accelerated odd trajectories are imposed explicitly through
admissibility rather than through the ambient definition of \(\mathbb N\).

Thus the occurrence of \(0\) in the symbolic and arithmetic constructions
does not enlarge the dynamical domain of the Collatz problem.


% ============================================================
\section{Division Counts and Division Words}
\label{sec:division-counts}
% ============================================================

% ------------------------------------------------------------
\subsection{Division Counts and Division Words}
% ------------------------------------------------------------

For a positive integer \(n\), let \(v_2(n)\) denote the largest
\(k\ge0\) such that \(2^k\mid n\). For \(x\in\mathbb N\), define the
\emph{division count}
\[
d(x):=v_2(3x+1),
\]
and define
\[
f(x):=\frac{3x+1}{2^{d(x)}}.
\]
These formulas define arithmetic functions on all of \(\mathbb N\).
Restricted to the odd positive integers, \(f\) is the fully accelerated
Collatz map.

Since \(2^{d(x)}\mid 3x+1\), the value \(f(x)\) is again a natural number.
By construction, removing the complete power of \(2\) from \(3x+1\) leaves
an odd integer. Hence
\[
f(x)\ \text{is odd for every }x\in\mathbb N.
\]
Moreover, if \(x\) is odd, then \(3x+1\) is even, and therefore
\[
d(x)\ge1.
\]

\begin{definition}[Division Word]
A \emph{division word} is a finite sequence
\[
\omega=(d_0,\dots,d_{m-1}),
\]
where \(m\ge0\) and each \(d_i\in\mathbb N\).
The case \(m=0\) is the empty division word, denoted by
\(\varnothing\).
\end{definition}

For \(x\in\mathbb N\) and \(m\ge0\), define the length-\(m\) division
word generated by \(x\) to be
\[
\omega_m(x)
:=
\bigl(
d(x),d(f(x)),\dots,d(f^{m-1}(x))
\bigr),
\]
with
\[
\omega_0(x)=\varnothing.
\]

\begin{definition}[Valid Division Word]
A division word
\[
\omega=(d_0,\dots,d_{m-1})
\]
is \emph{valid} if it is nonempty, equivalently if
\[
m\ge1.
\]
\end{definition}

\begin{definition}[Admissible Division Word]
A division word
\[
\omega=(d_0,\dots,d_{m-1})
\]
is \emph{admissible} if it is valid and
\[
d_i\ge1,
\qquad
0\le i<m.
\]
\end{definition}

\begin{definition}[Realization]
Let
\[
\omega=(d_0,\dots,d_{m-1})
\]
be a division word. A natural number \(x\in\mathbb N\) \emph{realizes}
\(\omega\) if
\[
\omega=\omega_m(x).
\]
Equivalently,
\[
d_i=v_2\!\left(3f^i(x)+1\right),
\qquad
0\le i<m.
\]
\end{definition}

A division word is called \emph{realizable} if it is realized by at least
one natural number.

\begin{lemma}[Generated Division Words]
\label{lem:generated-division-words}
For every \(x\in\mathbb N\) and every \(m\ge0\), the generated word
\(\omega_m(x)\) is a division word of length \(m\) and is realized by
\(x\).

If \(x\) is odd and \(m\ge1\), then \(\omega_m(x)\) is admissible.
\end{lemma}

\begin{proof}
For every \(i\ge0\),
\[
d(f^i(x))\in\mathbb N,
\]
so \(\omega_m(x)\) is a division word. Its length is \(m\), and realization
by \(x\) follows directly from its definition.

Now suppose that \(x\) is odd and \(m\ge1\). Then
\[
d(x)\ge1.
\]
Every subsequent iterate \(f^i(x)\) is odd, because \(f\) removes the
complete power of \(2\) from the corresponding numerator. Hence
\[
d(f^i(x))\ge1,
\qquad
0\le i<m.
\]
Thus every entry of the nonempty word \(\omega_m(x)\) is positive, and
\(\omega_m(x)\) is admissible.
\end{proof}

\begin{remark}
Realization and admissibility are distinct notions. Admissibility is an
intrinsic positivity condition on a nonempty symbolic word, whereas
realization asserts that the word actually occurs as an initial segment
generated by iteration of \(f\).

In particular, no converse to the preceding lemma is being assumed here.
The arithmetic characterization of realizable division words will be
established later.
\end{remark}

% ------------------------------------------------------------
\subsection{Word Length and Total Division Count}
% ------------------------------------------------------------

\begin{definition}[Length]
Let
\[
\omega=(d_0,\ldots,d_{m-1})
\]
be a division word. Its \emph{length} is
\[
|\omega|:=m.
\]
\end{definition}

\begin{definition}[Total Division Count]
Let
\[
\omega=(d_0,\ldots,d_{m-1})
\]
be a division word. Its \emph{total division count} is
\[
D(\omega):=\sum_{i=0}^{m-1}d_i.
\]
For the empty word,
\[
D(\varnothing)=0.
\]
\end{definition}

\begin{lemma}[Length Bound]
Every admissible division word satisfies
\[
|\omega|\le D(\omega).
\]
\end{lemma}

\begin{proof}
If
\[
\omega=(d_0,\ldots,d_{m-1})
\]
is admissible, then \(d_i\ge1\) for every \(0\le i<m\). Hence
\[
D(\omega)
=
\sum_{i=0}^{m-1}d_i
\ge
\sum_{i=0}^{m-1}1
=
m
=
|\omega|.
\]
\end{proof}

\begin{remark}
The quantities
\[
|\omega|
\qquad\text{and}\qquad
D(\omega)
\]
record different aspects of the symbolic trajectory. The length
\(|\omega|\) is the number of applications of \(f\) encoded by the word,
whereas \(D(\omega)\) is the cumulative exponent of \(2\) removed across
those steps.

Their interaction will determine the linear coefficient
\[
\frac{3^{|\omega|}}{2^{D(\omega)}}
\]
of the affine representation developed below.
\end{remark}


% ============================================================
\section{Affine Representation of Division Words}
\label{sec:affine-representation}
% ============================================================

% ------------------------------------------------------------
\subsection{Formal Symbolic Action and Affine Representation}
% ------------------------------------------------------------

\begin{definition}[Symbolic Step]
For a prescribed division count \(d\in\mathbb N\), define the
\emph{symbolic step}
\[
S_d:\mathbb Q\to\mathbb Q,
\qquad
S_d(x):=\frac{3x+1}{2^d}.
\]
No divisibility or realizability condition is imposed. Thus \(S_d\) is
defined for every \(d\in\mathbb N\) and every \(x\in\mathbb Q\).
\end{definition}

\begin{definition}[Symbolic Action of a Division Word]
Let
\[
\omega=(d_0,\dots,d_{m-1})
\]
be a division word. Its \emph{symbolic action} is the composition
\[
S_\omega
:=
S_{d_{m-1}}\circ\cdots\circ S_{d_0}.
\]
For the empty word,
\[
S_{\varnothing}(x):=x.
\]
Equivalently, if
\[
\omega=(d_0,\eta),
\]
where \(\eta\) denotes the remaining tail, then
\[
S_\omega(x)=S_\eta(S_{d_0}(x)).
\]
\end{definition}

\begin{definition}[Affine Constant]
The affine constant \(c(\omega)\in\mathbb N\) of a division word is
defined recursively by
\[
c(\varnothing)=0,
\]
and, for
\[
\omega=(d_0,\dots,d_{m-1}),\qquad m\ge1,
\]
by
\[
c(\omega)
=
3^{m-1}
+
2^{d_0}c(d_1,\dots,d_{m-1}).
\]
\end{definition}

\begin{definition}[Formal Affine Map]
For every division word \(\omega\), define
\[
F_\omega:\mathbb Q\to\mathbb Q
\]
by
\[
F_\omega(x)
=
\frac{3^{|\omega|}x+c(\omega)}
     {2^{D(\omega)}}.
\]
The rational domain is used because an arbitrary formal division word need
not prescribe powers of \(2\) that divide the corresponding numerators.

For the empty word,
\[
F_{\varnothing}(x)=x,
\]
since
\[
|\varnothing|=D(\varnothing)=c(\varnothing)=0.
\]
\end{definition}

\begin{theorem}[Affine Representation]
For every division word \(\omega\) and every \(x\in\mathbb Q\),
\[
S_\omega(x)=F_\omega(x).
\]
Equivalently, symbolic composition of the prescribed steps has the
closed affine form
\[
S_\omega(x)
=
\frac{3^{|\omega|}x+c(\omega)}
     {2^{D(\omega)}}.
\]
\end{theorem}

\begin{proof}
We argue by induction on the word.

For the empty word,
\[
S_{\varnothing}(x)=x=F_{\varnothing}(x).
\]

Now let
\[
\omega=(d_0,\eta),
\]
where \(\eta\) is the tail of the word. By the induction hypothesis,
\[
S_\eta(y)
=
\frac{3^{|\eta|}y+c(\eta)}
     {2^{D(\eta)}}
\]
for every \(y\in\mathbb Q\). Therefore
\[
\begin{aligned}
S_\omega(x)
&=
S_\eta\!\left(\frac{3x+1}{2^{d_0}}\right)\\
&=
\frac{
3^{|\eta|}
\left(\frac{3x+1}{2^{d_0}}\right)
+c(\eta)}
{2^{D(\eta)}}\\
&=
\frac{
3^{|\eta|+1}x
+
3^{|\eta|}
+
2^{d_0}c(\eta)}
{2^{d_0+D(\eta)}}.
\end{aligned}
\]
Since
\[
|\omega|=|\eta|+1,
\qquad
D(\omega)=d_0+D(\eta),
\]
and
\[
c(\omega)
=
3^{|\eta|}
+
2^{d_0}c(\eta),
\]
we obtain
\[
S_\omega(x)
=
\frac{3^{|\omega|}x+c(\omega)}
     {2^{D(\omega)}}
=
F_\omega(x).
\]
\end{proof}

\begin{corollary}[Affine Representation of a Realized Word]
If \(x\in\mathbb N\) realizes a division word \(\omega\), then
\[
F_\omega(x)=f^{|\omega|}(x).
\]
\end{corollary}

\begin{proof}
Write
\[
\omega=(d_0,\dots,d_{m-1}).
\]
Since \(x\) realizes \(\omega\),
\[
d_i
=
v_2\!\left(3f^i(x)+1\right),
\qquad
0\le i<m.
\]
Consequently, at each stage the prescribed symbolic step agrees with the
corresponding application of the arithmetic map \(f\):
\[
S_{d_i}(f^i(x))=f^{i+1}(x).
\]
Iterating through the word gives
\[
S_\omega(x)=f^m(x).
\]
The affine representation theorem then yields
\[
F_\omega(x)=S_\omega(x)=f^{|\omega|}(x).
\]
\end{proof}

\begin{remark}[Formal Integrality and Exact Realization]
The denominator in the affine representation is
\[
2^{D(\omega)}.
\]
Thus, for \(x\in\mathbb N\), the formal affine value \(F_\omega(x)\) is an
integer precisely when
\[
3^{|\omega|}x+c(\omega)
\equiv
0
\pmod{2^{D(\omega)}}.
\]

This condition expresses integrality of the formal affine expression only.
It does not by itself assert that \(x\) realizes \(\omega\). Exact
realizability requires the prescribed powers of \(2\) to be the exact
successive \(2\)-adic valuations, a stronger condition developed below.
\end{remark}

% ------------------------------------------------------------
\subsection{Concatenation and Prefix Structure}
% ------------------------------------------------------------

Let \(\omega\eta\) denote the concatenation obtained by appending the
division word \(\eta\) to the division word \(\omega\). Then
\[
|\omega\eta|=|\omega|+|\eta|
\]
and
\[
D(\omega\eta)=D(\omega)+D(\eta).
\]

\begin{proposition}[Concatenation of Symbolic and Affine Actions]
For all division words \(\omega\) and \(\eta\),
\[
S_{\omega\eta}
=
S_\eta\circ S_\omega
\]
and consequently
\[
F_{\omega\eta}
=
F_\eta\circ F_\omega.
\]
\end{proposition}

\begin{proof}
The first identity follows directly from the definition of symbolic action:
the entries of \(\omega\) are applied first, followed by the entries of
\(\eta\).

By the affine representation theorem,
\[
S_\omega=F_\omega,
\qquad
S_\eta=F_\eta,
\qquad
S_{\omega\eta}=F_{\omega\eta}.
\]
Therefore
\[
F_{\omega\eta}
=
S_{\omega\eta}
=
S_\eta\circ S_\omega
=
F_\eta\circ F_\omega.
\]
\end{proof}

\begin{corollary}[Affine Constant under Concatenation]
For all division words \(\omega\) and \(\eta\),
\[
c(\omega\eta)
=
3^{|\eta|}c(\omega)
+
2^{D(\omega)}c(\eta).
\]
\end{corollary}

\begin{proof}
By the preceding proposition,
\[
F_{\omega\eta}(x)=F_\eta(F_\omega(x)).
\]
Using the affine representations of \(F_\omega\) and \(F_\eta\),
\[
\begin{aligned}
F_\eta(F_\omega(x))
&=
\frac{
3^{|\eta|}
\left(
\frac{3^{|\omega|}x+c(\omega)}
     {2^{D(\omega)}}
\right)
+c(\eta)}
{2^{D(\eta)}}\\
&=
\frac{
3^{|\omega|+|\eta|}x
+
3^{|\eta|}c(\omega)
+
2^{D(\omega)}c(\eta)}
{2^{D(\omega)+D(\eta)}}.
\end{aligned}
\]
On the other hand,
\[
F_{\omega\eta}(x)
=
\frac{
3^{|\omega|+|\eta|}x+c(\omega\eta)}
{2^{D(\omega)+D(\eta)}}.
\]
Comparing the constant terms gives
\[
c(\omega\eta)
=
3^{|\eta|}c(\omega)
+
2^{D(\omega)}c(\eta).
\]
\end{proof}

For a division word \(\omega\), let
\[
R_\omega
:=
\{x\in\mathbb N:x\text{ realizes }\omega\}
\]
denote its realization set.

\begin{proposition}[Prefix Nesting]
For all division words \(\omega\) and \(\eta\),
\[
R_{\omega\eta}\subseteq R_\omega.
\]
Equivalently, if \(\omega\) is a prefix of a division word \(\nu\), then
\[
R_\nu\subseteq R_\omega.
\]
\end{proposition}

\begin{proof}
Suppose \(x\in R_{\omega\eta}\). Then the initial
\(|\omega|+|\eta|\) division counts generated by \(x\) are precisely the
entries of \(\omega\eta\). Their first \(|\omega|\) entries are therefore
the entries of \(\omega\). Hence \(x\) realizes \(\omega\), and so
\[
x\in R_\omega.
\]
Thus
\[
R_{\omega\eta}\subseteq R_\omega.
\]
The equivalent prefix formulation follows because any extension of
\(\omega\) can be written as \(\omega\eta\) for some division word
\(\eta\).
\end{proof}

\begin{remark}
Concatenation therefore has parallel algebraic and dynamical meanings:
\[
\omega\longmapsto\omega\eta
\]
corresponds algebraically to composition of affine maps,
\[
F_{\omega\eta}=F_\eta\circ F_\omega,
\]
and dynamically to refinement of realization sets,
\[
R_{\omega\eta}\subseteq R_\omega.
\]
The latter nesting principle will become the arithmetic nesting of exact
dyadic realization classes once realizability is characterized
congruentially.
\end{remark}

% ------------------------------------------------------------
\subsection{Explicit Form of the Affine Constant}
% ------------------------------------------------------------

The recursive definition of \(c(\omega)\) is convenient for constructing
the affine map, but it obscures the contribution of the individual
additive terms introduced along the symbolic trajectory. We now make those
contributions explicit.

Let
\[
\omega=(d_0,\dots,d_{m-1})
\]
be a division word. For \(0\le j\le m\), define the \emph{prefix division
count}
\[
D_j(\omega)
:=
\sum_{i=0}^{j-1}d_i.
\]
Thus
\[
D_0(\omega)=0,
\qquad
D_m(\omega)=D(\omega).
\]

\begin{theorem}[Explicit Form of the Affine Constant]
\label{thm:explicit-affine-constant}
For every division word
\[
\omega=(d_0,\dots,d_{m-1}),
\]
the affine constant is
\[
c(\omega)
=
\sum_{j=0}^{m-1}
3^{m-1-j}2^{D_j(\omega)}.
\]
Equivalently,
\[
c(\omega)
=
\sum_{j=0}^{m-1}
3^{m-1-j}
2^{\sum_{i=0}^{j-1}d_i},
\]
where the sum in the exponent is \(0\) when \(j=0\).
\end{theorem}

\begin{proof}
We argue by induction on the length \(m\).

For \(m=0\), the word is empty. Both the defining value
\[
c(\varnothing)=0
\]
and the displayed sum are zero.

Now suppose \(m\ge1\), and write
\[
\omega=(d_0,\eta),
\qquad
\eta=(d_1,\dots,d_{m-1}).
\]
By the recursive definition of the affine constant,
\[
c(\omega)
=
3^{m-1}
+
2^{d_0}c(\eta).
\]
Applying the induction hypothesis to the tail \(\eta\) gives
\[
c(\eta)
=
\sum_{k=0}^{m-2}
3^{m-2-k}
2^{\sum_{i=1}^{k}d_i},
\]
where the exponent sum is empty when \(k=0\). Hence
\[
c(\omega)
=
3^{m-1}
+
\sum_{k=0}^{m-2}
3^{m-2-k}
2^{d_0+\sum_{i=1}^{k}d_i}.
\]
Set \(j=k+1\). Then
\[
m-2-k=m-1-j
\]
and
\[
d_0+\sum_{i=1}^{k}d_i
=
\sum_{i=0}^{j-1}d_i
=
D_j(\omega).
\]
The initial term is
\[
3^{m-1}
=
3^{m-1}2^{D_0(\omega)},
\]
so adjoining it as the \(j=0\) term yields
\[
c(\omega)
=
\sum_{j=0}^{m-1}
3^{m-1-j}2^{D_j(\omega)}.
\]
\end{proof}

\begin{corollary}
For every division word \(\omega=(d_0,\dots,d_{m-1})\),
\[
F_\omega(x)
=
\frac{
3^m x
+
\displaystyle\sum_{j=0}^{m-1}
3^{m-1-j}2^{D_j(\omega)}
}
{2^{D(\omega)}}.
\]
\end{corollary}

\begin{remark}
The explicit formula records the accumulated contribution of each
additive \(+1\) introduced during symbolic iteration. The contribution
created at step \(j\) is multiplied by
\[
3^{m-1-j}
\]
during the remaining symbolic steps, while the factor
\[
2^{D_j(\omega)}
\]
records all prescribed divisions that occurred before that contribution
was introduced.

Thus the affine constant retains the time ordering of the division word,
not merely its total division count.
\end{remark}

% ------------------------------------------------------------
\subsection{Dyadic-Ternary Balance}
% ------------------------------------------------------------

The affine representation
\[
F_\omega(x)
=
\frac{3^{|\omega|}x+c(\omega)}
     {2^{D(\omega)}}
\]
contains two cumulative multiplicative factors: the ternary factor
\(3^{|\omega|}\), contributed by the \(3x+1\) term at each symbolic step,
and the dyadic factor \(2^{D(\omega)}\), arising from the prescribed
divisions.

Their balance is governed by a natural integer threshold.

\begin{definition}[Dyadic Threshold]
For \(m\in\mathbb N\), let \(B(m)\) be the least natural number \(k\)
such that
\[
3^m<2^k.
\]
Thus
\[
B(m)
=
\min\{k\in\mathbb N:3^m<2^k\}.
\]
\end{definition}

Since \(0\in\mathbb N\), the definition includes \(m=0\). In particular,
\[
B(0)=1,
\]
since \(3^0=1<2\), while \(1\nless 1\).

\begin{lemma}[Threshold Characterization]
For all \(m,k\in\mathbb N\),
\[
B(m)\le k
\quad\Longleftrightarrow\quad
3^m<2^k.
\]
\end{lemma}

\begin{proof}
By definition, \(B(m)\) is the least exponent for which
\[
3^m<2^{B(m)}.
\]
If \(B(m)\le k\), monotonicity of \(2^k\) gives
\[
3^m<2^{B(m)}\le2^k.
\]
Conversely, if \(3^m<2^k\), then \(k\) belongs to the defining set for
\(B(m)\), and minimality gives
\[
B(m)\le k.
\]
\end{proof}

\begin{proposition}[Logarithmic Form of the Threshold]
For every \(m\in\mathbb N\),
\[
B(m)
=
\lfloor m\log_2 3\rfloor+1.
\]
\end{proposition}

\begin{proof}
The inequality
\[
3^m<2^k
\]
is equivalent to
\[
m\log_2 3<k.
\]
Hence \(B(m)\) is the least integer strictly greater than
\(m\log_2 3\), namely
\[
\lfloor m\log_2 3\rfloor+1.
\]
\end{proof}

\begin{remark}
The sequence \(B(m)\) is OEIS sequence \OEIS{A020914}.
The OEIS identification is descriptive only; the argument below uses the
intrinsic threshold characterization.
\end{remark}

\begin{definition}[Affine Slope]
For a division word \(\omega\), define its \emph{affine slope} by
\[
\lambda_\omega
:=
\frac{3^{|\omega|}}{2^{D(\omega)}}.
\]
Thus
\[
F_\omega(x)
=
\lambda_\omega x
+
\frac{c(\omega)}{2^{D(\omega)}}.
\]
\end{definition}

\begin{theorem}[Dyadic Threshold and Affine Slope]
For every division word \(\omega\),
\[
\lambda_\omega<1
\quad\Longleftrightarrow\quad
B(|\omega|)\le D(\omega).
\]
Equivalently,
\[
\lambda_\omega<1
\quad\Longleftrightarrow\quad
3^{|\omega|}<2^{D(\omega)}.
\]
\end{theorem}

\begin{proof}
Since \(D(\omega)\in\mathbb N\), we have
\[
2^{D(\omega)}\ge1,
\]
and hence the denominator of \(\lambda_\omega\) is positive. Therefore
\[
\lambda_\omega<1
\quad\Longleftrightarrow\quad
3^{|\omega|}<2^{D(\omega)}.
\]
Applying the threshold characterization with
\[
m=|\omega|,
\qquad
k=D(\omega)
\]
gives
\[
3^{|\omega|}<2^{D(\omega)}
\quad\Longleftrightarrow\quad
B(|\omega|)\le D(\omega).
\]
\end{proof}

\begin{corollary}[Strict Slope Dichotomy]
If \(\omega\) is nonempty, then
\[
\lambda_\omega<1
\qquad\text{or}\qquad
\lambda_\omega>1.
\]
In particular,
\[
\lambda_\omega\ne1.
\]
\end{corollary}

\begin{proof}
Equality would require
\[
3^{|\omega|}=2^{D(\omega)}.
\]
If \(D(\omega)>0\), the right-hand side is even while the left-hand side
is odd. If \(D(\omega)=0\), the right-hand side is \(1\), whereas
\(3^{|\omega|}>1\) because \(\omega\) is nonempty. Thus equality is
impossible.
\end{proof}

\begin{corollary}[Positivity of the Affine Constant]
If \(\omega\) is nonempty, then
\[
c(\omega)>0.
\]
\end{corollary}

\begin{proof}
Let \(|\omega|=m\ge1\). By the explicit formula,
\[
c(\omega)
=
\sum_{j=0}^{m-1}
3^{m-1-j}2^{D_j(\omega)}.
\]
Every summand is nonnegative, and the \(j=0\) term is
\[
3^{m-1}>0.
\]
Hence \(c(\omega)>0\).
\end{proof}

\begin{remark}[Slope versus Descent]
The inequality
\[
\lambda_\omega<1
\]
is a statement about the linear coefficient of \(F_\omega\); it is not by
itself a statement that the encoded trajectory descends below its starting
value.

Indeed,
\[
F_\omega(x)-x
=
\frac{
\bigl(3^{|\omega|}-2^{D(\omega)}\bigr)x+c(\omega)}
{2^{D(\omega)}}.
\]
When
\[
3^{|\omega|}<2^{D(\omega)},
\]
the linear contribution is negative, but the affine constant
\(c(\omega)>0\) offsets part of that decrease. Thus the slope determines
the eventual direction of the affine map, while actual descent also
depends on the starting value and the affine constant.

The relation between subunit slope and genuine descent on realized
trajectories will be established later.
\end{remark}

% ------------------------------------------------------------
\subsection{Formal Integrality versus Exact Realization}
% ------------------------------------------------------------

Consider the division word
\[
\omega=(1,1,1,4).
\]
It has length
\[
|\omega|=4
\]
and total division count
\[
D(\omega)=1+1+1+4=7.
\]
Since
\[
B(4)=7,
\]
this word lies exactly at the dyadic threshold:
\[
3^4=81<128=2^7,
\]
and therefore
\[
\lambda_\omega=\frac{81}{128}<1.
\]

The prefix division counts are
\[
D_0(\omega)=0,\qquad
D_1(\omega)=1,\qquad
D_2(\omega)=2,\qquad
D_3(\omega)=3.
\]
Hence the explicit formula for the affine constant gives
\[
\begin{aligned}
c(\omega)
&=
3^3 2^0
+
3^2 2^1
+
3^1 2^2
+
3^0 2^3\\
&=
27+18+12+8\\
&=
65.
\end{aligned}
\]
Thus
\[
F_\omega(x)
=
\frac{81x+65}{128}.
\]

For \(F_\omega(x)\) to be integral, it is necessary and sufficient that
\[
81x+65\equiv0\pmod{128}.
\]
Since
\[
81^{-1}\equiv49\pmod{128},
\]
this is equivalent to
\[
x\equiv15\pmod{128}.
\]

This residue class therefore describes exactly the natural numbers for which
the \emph{formal} affine expression associated with \(\omega\) is integral.
It does not, however, follow that every such number realizes \(\omega\).

Indeed, the least positive representative \(x=15\) satisfies the
integrality congruence, but its fully accelerated trajectory begins
\[
15\longmapsto23\longmapsto35\longmapsto53\longmapsto5.
\]
The first three division counts are \(1,1,1\), but at the fourth step
\[
3\cdot53+1
=
160
=
2^5\cdot5.
\]
Thus the fourth exact valuation is \(5\), whereas \(\omega\) prescribes
\(4\).

The formal symbolic step prescribed by \(\omega\) would instead give
\[
\frac{3\cdot53+1}{2^4}
=
10.
\]
Accordingly,
\[
F_\omega(15)=10,
\]
while the actual fourth fully accelerated iterate is
\[
f^4(15)=5.
\]
The formal affine expression is integral, but \(15\) does not realize
\(\omega\).

If \(x\) did realize \(\omega\), then
\[
F_\omega(x)=f^4(x)
\]
would be odd. Hence exact realization necessarily imposes the additional
parity condition
\[
\frac{81x+65}{128}\equiv1\pmod2,
\]
or equivalently
\[
81x+65\equiv128\pmod{256}.
\]
Since
\[
81^{-1}\equiv177\pmod{256},
\]
this strengthened congruence is equivalent to
\[
x\equiv143\pmod{256}.
\]

Taking the least positive representative \(x=143\), we find
\[
143\longmapsto215\longmapsto323\longmapsto485\longmapsto91.
\]
Indeed,
\[
\begin{aligned}
3\cdot143+1 &= 2\cdot215,\\
3\cdot215+1 &= 2\cdot323,\\
3\cdot323+1 &= 2\cdot485,\\
3\cdot485+1 &= 2^4\cdot91,
\end{aligned}
\]
so the successive exact \(2\)-adic valuations are
\[
(1,1,1,4).
\]
Thus \(143\) genuinely realizes \(\omega\).

\begin{remark}[Interpretation]
The example exhibits two distinct arithmetic conditions.

The congruence
\[
3^{|\omega|}x+c(\omega)
\equiv0
\pmod{2^{D(\omega)}}
\]
expresses only integrality of the formal affine map. For the present word,
it gives the class
\[
x\equiv15\pmod{128},
\]
which contains values such as \(15\) that do not realize \(\omega\).

Realization necessarily requires the quotient obtained after removing the
prescribed total power \(2^{D(\omega)}\) to be odd. In this example that
adds one dyadic bit and refines the integrality class to
\[
x\equiv143\pmod{256}.
\]

The direct calculation above shows that \(143\) realizes the word. It does
not yet establish that the strengthened congruence characterizes all
realizations of an arbitrary admissible division word. The following
development proves precisely that stronger statement.
\end{remark}


% ============================================================
\section{Exact Realizability and Residue Structure}
\label{sec:exact-realizability-residue}
% ============================================================

The preceding section associated a formal affine map with every division
word, independently of whether that word is realized under iteration of the
arithmetic map \(f\). We now determine the additional arithmetic conditions
that distinguish formal symbolic integrality from exact realization.

% ------------------------------------------------------------
\subsection{Affine Integrality and Prefix Propagation}
% ------------------------------------------------------------

\begin{theorem}[Natural Affine Identity for a Realized Word]
\label{thm:natural-affine-identity}

Let
\[
\omega=(d_0,\dots,d_{m-1})
\]
be a division word, and let \(x\in\mathbb N\) realize \(\omega\). Then
\[
3^m x+c(\omega)
=
2^{D(\omega)}f^m(x).
\]
\end{theorem}

\begin{proof}
Since \(x\) realizes \(\omega\), the affine representation of a realized
word gives
\[
F_\omega(x)=f^m(x).
\]
By definition,
\[
F_\omega(x)
=
\frac{3^m x+c(\omega)}{2^{D(\omega)}}.
\]
Hence
\[
\frac{3^m x+c(\omega)}{2^{D(\omega)}}
=
f^m(x).
\]
Multiplying by \(2^{D(\omega)}\) yields
\[
3^m x+c(\omega)
=
2^{D(\omega)}f^m(x).
\]
\end{proof}

\begin{corollary}[Formal Integrality of a Realized Word]
\label{cor:realized-formal-integrality}

If \(x\in\mathbb N\) realizes a division word \(\omega\), then
\[
3^{|\omega|}x+c(\omega)
\equiv
0
\pmod{2^{D(\omega)}}.
\]
\end{corollary}

\begin{proof}
By Theorem~\ref{thm:natural-affine-identity},
\[
3^{|\omega|}x+c(\omega)
=
2^{D(\omega)}f^{|\omega|}(x),
\]
so the affine numerator is divisible by \(2^{D(\omega)}\).
\end{proof}

\begin{remark}
The converse is not yet available. Divisibility of the affine numerator by
\(2^{D(\omega)}\) guarantees only that the formal affine endpoint is
integral. It does not assert that the prescribed powers of \(2\) are the
exact successive \(2\)-adic valuations.

The example of \(\omega=(1,1,1,4)\) showed this distinction explicitly:
\(x=15\) satisfies the formal integrality congruence but does not realize
the word.
\end{remark}

\begin{theorem}[Prefix Congruence Propagation]
\label{thm:prefix-congruence-propagation}

Let
\[
\omega=(d_0,\dots,d_{m-1})
\]
be a division word, let \(x\in\mathbb N\), and suppose
\[
3^m x+c(\omega)
\equiv
0
\pmod{2^{D(\omega)}}.
\]
For \(0\le j\le m\), let
\[
\omega_{\le j}
=
(d_0,\dots,d_{j-1})
\]
denote the prefix of length \(j\). Then
\[
3^j x+c(\omega_{\le j})
\equiv
0
\pmod{2^{D(\omega_{\le j})}}.
\]
Thus formal integrality of the full affine numerator forces the
corresponding formal integrality condition at every prefix.
\end{theorem}

\begin{proof}
Fix \(j\) with \(0\le j\le m\), and write
\[
\omega=\omega_{\le j}\eta,
\]
where \(\eta\) is the remaining suffix.

The concatenation formula for the affine constant gives
\[
c(\omega)
=
3^{|\eta|}c(\omega_{\le j})
+
2^{D(\omega_{\le j})}c(\eta),
\]
and
\[
m=j+|\eta|.
\]
Therefore
\[
3^m x+c(\omega)
=
3^{|\eta|}
\bigl(3^j x+c(\omega_{\le j})\bigr)
+
2^{D(\omega_{\le j})}c(\eta).
\]

Moreover,
\[
D(\omega)
=
D(\omega_{\le j})+D(\eta),
\]
so
\[
2^{D(\omega_{\le j})}
\mid
2^{D(\omega)}.
\]
The assumed congruence modulo \(2^{D(\omega)}\) therefore implies
\[
3^m x+c(\omega)
\equiv
0
\pmod{2^{D(\omega_{\le j})}}.
\]
The second term in the preceding decomposition is already divisible by
\(2^{D(\omega_{\le j})}\), hence
\[
3^{|\eta|}
\bigl(3^j x+c(\omega_{\le j})\bigr)
\equiv
0
\pmod{2^{D(\omega_{\le j})}}.
\]

Since \(3^{|\eta|}\) is odd, it is coprime to every power of \(2\) and is
therefore invertible modulo \(2^{D(\omega_{\le j})}\). Cancelling this
factor gives
\[
3^j x+c(\omega_{\le j})
\equiv
0
\pmod{2^{D(\omega_{\le j})}}.
\]
Since \(j\) was arbitrary, the result holds for every prefix.
\end{proof}

\begin{remark}[Endpoint Integrality versus Stepwise Integrality]
The prefix theorem shows that the global affine congruence is not merely an
endpoint condition. If
\[
3^{|\omega|}x+c(\omega)
\equiv
0
\pmod{2^{D(\omega)}},
\]
then for every prefix \(\omega_{\le j}\),
\[
F_{\omega_{\le j}}(x)
=
\frac{
3^j x+c(\omega_{\le j})
}{
2^{D(\omega_{\le j})}
}
\]
is a natural number.

Thus the prescribed symbolic divisions are arithmetically compatible at
every prefix stage. What remains unresolved is \emph{exactness}: the
quotient after a prescribed division may still contain an additional
factor of \(2\).

Exact realization therefore requires one further parity bit. The next
subsection identifies that local exactness condition and then lifts it to
the complete division word.
\end{remark}

% ------------------------------------------------------------
\subsection{Exact Valuation and the Additional Parity Bit}
% ------------------------------------------------------------

Formal integrality guarantees that the prescribed power of \(2\) divides
the relevant numerator. Exact realization requires more: after that power
of \(2\) has been removed, the remaining quotient must be odd.

The local arithmetic condition is elementary but fundamental.

\begin{lemma}[Exact Valuation Congruence]
\label{lem:exact-valuation-congruence}
Let \(x,d\in\mathbb N\). Then
\[
v_2(3x+1)=d
\]
if and only if
\[
3x+1
\equiv
2^d
\pmod{2^{d+1}}.
\]
Equivalently,
\[
3x+1=2^d q
\]
for some odd \(q\in\mathbb N\).
\end{lemma}

\begin{proof}
If
\[
v_2(3x+1)=d,
\]
then by definition
\[
3x+1=2^d q
\]
for some odd natural number \(q\). Writing
\[
q=2k+1
\]
gives
\[
3x+1
=
2^d(2k+1)
=
2^{d+1}k+2^d,
\]
and hence
\[
3x+1
\equiv
2^d
\pmod{2^{d+1}}.
\]

Conversely, suppose
\[
3x+1
\equiv
2^d
\pmod{2^{d+1}}.
\]
Then
\[
3x+1
=
2^{d+1}k+2^d
=
2^d(2k+1)
\]
for some \(k\in\mathbb N\). Since \(2k+1\) is odd, exactly \(d\) factors
of \(2\) divide \(3x+1\). Therefore
\[
v_2(3x+1)=d.
\]
\end{proof}

\begin{remark}
The distinction between divisibility and exact valuation is therefore the
distinction between
\[
3x+1\equiv0\pmod{2^d}
\]
and
\[
3x+1\equiv2^d\pmod{2^{d+1}}.
\]

The first says only that the prescribed division by \(2^d\) is possible.
The second says that the quotient after that division is odd, and therefore
that \(d\) is the exact \(2\)-adic valuation.
\end{remark}


% ------------------------------------------------------------
\subsection{The Exact Realization Criterion}
% ------------------------------------------------------------

The preceding lemma gives a local criterion for one exact symbolic step.
The next theorem shows that, for an admissible division word, all of the
local exact-valuation conditions are encoded simultaneously by one global
affine congruence.

\begin{theorem}[Exact Realization Criterion]
\label{thm:exact-realization-criterion}

Let
\[
\omega=(d_0,\dots,d_{m-1})
\]
be an admissible division word, and let \(x\in\mathbb N\). Then \(x\)
realizes \(\omega\) if and only if
\[
3^m x+c(\omega)
\equiv
2^{D(\omega)}
\pmod{2^{D(\omega)+1}}.
\]
\end{theorem}

\begin{proof}
Suppose first that \(x\) realizes \(\omega\). Since \(\omega\) is
admissible, it is nonempty and its first division count satisfies
\(d_0\ge1\). Since \(x\) realizes \(\omega\),
\[
d_0=v_2(3x+1).
\]
Hence
\[
v_2(3x+1)\ge1,
\]
so \(x\) is odd. Since \(x\in\mathbb N\), it follows that \(x>0\).
Thus the realized trajectory is genuinely a fully accelerated Collatz
trajectory.

By the natural affine identity,
\[
3^m x+c(\omega)
=
2^{D(\omega)}f^m(x).
\]
Since \(m\ge1\), the accelerated endpoint \(f^m(x)\) is odd. Hence
\[
f^m(x)\equiv1\pmod2.
\]
Multiplying this congruence by \(2^{D(\omega)}\) gives
\[
3^m x+c(\omega)
\equiv
2^{D(\omega)}
\pmod{2^{D(\omega)+1}}.
\]

For the converse, we prove a slightly stronger statement by induction on
the word length.

Let
\[
\eta=(e_0,\dots,e_{n-1})
\]
be any finite division word whose entries are positive, allowing also the
empty word when \(n=0\). We claim that if
\[
3^n x+c(\eta)
\equiv
2^{D(\eta)}
\pmod{2^{D(\eta)+1}},
\]
then \(x\) is odd and realizes \(\eta\).

If \(n=0\), then
\[
D(\eta)=0,
\qquad
c(\eta)=0,
\]
and the exact congruence becomes
\[
x\equiv1\pmod2.
\]
Thus \(x\) is odd. The empty word is realized trivially, so the claim
holds.

Now suppose \(n\ge1\), and write
\[
\eta=(e_0,\eta'),
\qquad
\eta'=(e_1,\dots,e_{n-1}).
\]
The exact congruence implies, after reducing modulo \(2^{D(\eta)}\),
\[
3^n x+c(\eta)
\equiv
0
\pmod{2^{D(\eta)}}.
\]
By prefix congruence propagation, the one-term prefix \((e_0)\) therefore
satisfies
\[
3x+1
\equiv
0
\pmod{2^{e_0}}.
\]
Hence there exists \(y\in\mathbb N\) such that
\[
3x+1=2^{e_0}y.
\]

Using
\[
c(\eta)
=
3^{n-1}
+
2^{e_0}c(\eta')
\]
and
\[
D(\eta)
=
e_0+D(\eta'),
\]
we obtain
\[
\begin{aligned}
3^n x+c(\eta)
&=
3^{n-1}(3x+1)
+
2^{e_0}c(\eta')\\
&=
2^{e_0}
\left(
3^{n-1}y+c(\eta')
\right).
\end{aligned}
\]

Substituting this factorization into the exact congruence gives
\[
2^{e_0}
\left(
3^{n-1}y+c(\eta')
\right)
\equiv
2^{e_0}2^{D(\eta')}
\pmod{
2^{e_0}2^{D(\eta')+1}
}.
\]
Since the modulus itself contains the factor \(2^{e_0}\), this congruence
is equivalent to divisibility by
\[
2^{e_0}2^{D(\eta')+1}
\]
of the difference of the two sides. Cancelling the common positive factor
\(2^{e_0}\) in that divisibility relation yields
\[
3^{n-1}y+c(\eta')
\equiv
2^{D(\eta')}
\pmod{2^{D(\eta')+1}}.
\]

Every entry of \(\eta'\) is positive, so the induction hypothesis applies.
Therefore \(y\) is odd and realizes \(\eta'\).

Since
\[
3x+1=2^{e_0}y
\]
with \(y\) odd, Lemma~\ref{lem:exact-valuation-congruence} gives
\[
v_2(3x+1)=e_0.
\]
Consequently,
\[
f(x)=y.
\]
Moreover \(e_0\ge1\), so \(3x+1\) is even, which forces \(x\) to be odd.

Thus the first exact division count of \(x\) is \(e_0\), and its next
accelerated iterate \(y=f(x)\) realizes the tail \(\eta'\). Therefore
\(x\) realizes the complete word \(\eta\).

This proves the induction statement. Applying it to the admissible word
\(\omega\) proves the converse implication.
\end{proof}

\begin{remark}[The Additional Dyadic Bit]
The distinction between the two affine congruences is now exact.

The congruence
\[
3^{|\omega|}x+c(\omega)
\equiv
0
\pmod{2^{D(\omega)}}
\]
expresses formal integrality. It guarantees that all prescribed divisions
are arithmetically possible, including at every prefix, but does not require
those divisions to remove the complete power of \(2\).

By contrast,
\[
3^{|\omega|}x+c(\omega)
\equiv
2^{D(\omega)}
\pmod{2^{D(\omega)+1}}
\]
is equivalent, for admissible \(\omega\), to exact realization of the
entire division word.

The additional dyadic bit records that, after the prescribed factor
\(2^{D(\omega)}\) has been removed, the resulting endpoint is odd.
Remarkably, together with prefix integrality this single endpoint parity
condition forces exactness at every intermediate step.
\end{remark}


% ------------------------------------------------------------
\subsection{Exact Realization Residue Classes}
% ------------------------------------------------------------

The exact realization criterion converts realization into a single linear
congruence in the initial value. Because its coefficient is an odd power of
\(3\), that congruence determines one and only one dyadic residue class.

\begin{theorem}[Exact Realization Residue Class]
\label{thm:exact-realization-residue-class}

Let
\[
\omega=(d_0,\dots,d_{m-1})
\]
be an admissible division word. Then there exists a unique residue class
\[
r_\omega
\pmod{2^{D(\omega)+1}}
\]
such that
\[
x\text{ realizes }\omega
\quad\Longleftrightarrow\quad
x\equiv r_\omega
\pmod{2^{D(\omega)+1}}.
\]

Equivalently, there is a unique canonical representative
\[
0\le r_\omega<2^{D(\omega)+1}
\]
with this property.
\end{theorem}

\begin{proof}
By Theorem~\ref{thm:exact-realization-criterion}, \(x\) realizes
\(\omega\) if and only if
\[
3^m x+c(\omega)
\equiv
2^{D(\omega)}
\pmod{2^{D(\omega)+1}}.
\]

Since \(3^m\) is odd,
\[
\gcd\!\left(3^m,2^{D(\omega)+1}\right)=1.
\]
Hence multiplication by \(3^m\) is invertible modulo
\(2^{D(\omega)+1}\). Therefore the affine linear congruence
\[
3^m x+c(\omega)
\equiv
2^{D(\omega)}
\pmod{2^{D(\omega)+1}}
\]
has a unique solution class for \(x\) modulo \(2^{D(\omega)+1}\).

Denote its unique canonical representative by \(r_\omega\). Then
\[
x\text{ realizes }\omega
\quad\Longleftrightarrow\quad
x\equiv r_\omega
\pmod{2^{D(\omega)+1}},
\]
as required.
\end{proof}

\begin{corollary}[Admissible Division Words Are Realizable]
\label{cor:admissible-words-realizable}

Every admissible division word is realizable.
\end{corollary}

\begin{proof}
Let \(\omega\) be admissible. By
Theorem~\ref{thm:exact-realization-residue-class}, its exact realization
congruence has a canonical solution
\[
r_\omega
\pmod{2^{D(\omega)+1}}.
\]
Taking \(x=r_\omega\) gives
\[
x\equiv r_\omega
\pmod{2^{D(\omega)+1}},
\]
and hence \(x\) realizes \(\omega\).
\end{proof}

\begin{remark}
Admissibility was defined only as an intrinsic positivity condition on a
finite nonempty division word. We have now proved, rather than assumed,
that this condition is sufficient for finite realizability.

Since \(\omega\) is admissible, its first division count is positive.
Therefore any realizer of \(\omega\), and in particular \(r_\omega\), is
odd and positive.

\begin{center}
\fbox{%
\parbox{0.85\textwidth}{%
\centering
Every finite nonempty sequence of positive division counts
occurs along some fully accelerated Collatz trajectory.
}}
\end{center}

Moreover, its complete realization set is exactly one residue class
modulo
\[
2^{D(\omega)+1}.
\]

The next subsection studies how these exact realization classes behave when
division words are extended through their prefixes.
\end{remark}

% ------------------------------------------------------------
\subsection{Nested Structure of Realization Classes}
% ------------------------------------------------------------

For a division word \(\omega\), recall its realization set
\[
R_\omega
=
\{x\in\mathbb N : x\text{ realizes }\omega\}.
\]

If \(\omega\) is admissible, the preceding results show that \(R_\omega\)
is exactly one residue class modulo \(2^{D(\omega)+1}\):
\[
R_\omega
=
\left\{
x\in\mathbb N :
x\equiv r_\omega
\pmod{2^{D(\omega)+1}}
\right\},
\]
where
\[
0\le r_\omega<2^{D(\omega)+1}
\]
is uniquely determined.

We now show that the prefix structure of division words induces a
corresponding nested-or-disjoint structure on their realization sets.

\begin{lemma}[Common Realizations Force Prefix Comparability]
\label{lem:common-realization-prefix}

Let \(\omega\) and \(\eta\) be division words. If some \(x\in\mathbb N\)
realizes both \(\omega\) and \(\eta\), then one word is a prefix of the
other.

More precisely, if
\[
|\omega|\le|\eta|,
\]
then
\[
\omega=\eta_{\le|\omega|},
\]
where \(\eta_{\le|\omega|}\) denotes the prefix of \(\eta\) of length
\(|\omega|\).
\end{lemma}

\begin{proof}
Suppose that \(x\) realizes both \(\omega\) and \(\eta\), and assume
\[
|\omega|\le|\eta|.
\]
Since realization of a word implies realization of each of its prefixes,
\(x\) realizes the prefix
\[
\eta_{\le|\omega|}.
\]

Both \(\omega\) and \(\eta_{\le|\omega|}\) therefore have length
\(|\omega|\) and are realized by the same starting value \(x\). But a
fixed starting value determines a unique division word of each finite
length. Hence
\[
\omega=\eta_{\le|\omega|}.
\]

Interchanging the roles of \(\omega\) and \(\eta\) gives the opposite
case. Thus any two division words having a common realization are
prefix-comparable.
\end{proof}

\begin{theorem}[Nested-or-Disjoint Realization Classes]
\label{thm:realization-classes-nested-disjoint}

Let \(\omega\) and \(\eta\) be division words. Then at least one of
\[
R_\eta\subseteq R_\omega,
\qquad
R_\omega\subseteq R_\eta,
\qquad
R_\omega\cap R_\eta=\varnothing
\]
holds.

Equivalently, two realization sets are either nested or disjoint; they
cannot partially overlap.
\end{theorem}

\begin{proof}
If
\[
R_\omega\cap R_\eta=\varnothing,
\]
there is nothing to prove.

Otherwise choose
\[
x\in R_\omega\cap R_\eta.
\]
Then \(x\) realizes both \(\omega\) and \(\eta\), so by
Lemma~\ref{lem:common-realization-prefix} the two words are
prefix-comparable.

If \(\omega\) is a prefix of \(\eta\), then every natural number realizing
\(\eta\) also realizes \(\omega\). Hence
\[
R_\eta\subseteq R_\omega.
\]

Similarly, if \(\eta\) is a prefix of \(\omega\), then
\[
R_\omega\subseteq R_\eta.
\]

Thus any two realization sets are nested or disjoint.
\end{proof}

\begin{corollary}[Nested-or-Disjoint Exact Residue Classes]
\label{cor:exact-residue-classes-nested-disjoint}

Let \(\omega\) and \(\eta\) be admissible division words, with exact
realization residues
\[
r_\omega
\pmod{2^{D(\omega)+1}}
\qquad\text{and}\qquad
r_\eta
\pmod{2^{D(\eta)+1}}.
\]
Then their corresponding exact residue classes are either disjoint or one
is contained in the other.

Moreover, if \(\eta\) extends \(\omega\), then
\[
r_\eta
\equiv
r_\omega
\pmod{2^{D(\omega)+1}}.
\]
The analogous statement holds with \(\omega\) and \(\eta\) interchanged.
\end{corollary}

\begin{proof}
For admissible division words, the exact residue classes are precisely
their realization sets:
\[
x\equiv r_\omega\pmod{2^{D(\omega)+1}}
\quad\Longleftrightarrow\quad
x\in R_\omega,
\]
and similarly for \(\eta\).

The nested-or-disjoint conclusion therefore follows immediately from
Theorem~\ref{thm:realization-classes-nested-disjoint}.

Now suppose that \(\eta\) extends \(\omega\). The canonical representative
\(r_\eta\) belongs to its own exact realization class, and hence realizes
\(\eta\). It therefore realizes the prefix \(\omega\), so the residue
characterization of \(R_\omega\) gives
\[
r_\eta
\equiv
r_\omega
\pmod{2^{D(\omega)+1}}.
\]
\end{proof}

\begin{corollary}[Successive Prefix Refinement]
\label{cor:successive-prefix-refinement}

Let
\[
\omega=(d_0,\dots,d_{m-1})
\]
be admissible, and let
\[
\omega_j=(d_0,\dots,d_{j-1})
\]
for \(1\le j\le m\). Then, whenever \(1\le j<m\),
\[
R_{\omega_{j+1}}\subseteq R_{\omega_j}
\]
and
\[
r_{\omega_{j+1}}
\equiv
r_{\omega_j}
\pmod{2^{D(\omega_j)+1}}.
\]
Thus each additional exact valuation condition refines the realization
class of the preceding prefix.
\end{corollary}

\begin{proof}
Since \(\omega\) is admissible and \(1\le j<m\), both
\(\omega_j\) and \(\omega_{j+1}\) are nonempty admissible prefixes of
\(\omega\). Moreover, \(\omega_{j+1}\) extends \(\omega_j\).
Hence both conclusions follow from the preceding corollary.
\end{proof}

\begin{remark}
The realization classes therefore form a laminar family: extension of a
division word produces a nested refinement of its realization class,
whereas prefix-incomparable words have disjoint realization sets.

For admissible words these sets are not merely abstract dynamical classes;
each is a single dyadic residue class
\[
x\equiv r_\omega
\pmod{2^{D(\omega)+1}}.
\]
The next section develops the geometry and density of such dyadic residue
classes systematically by introducing them as cylinders.
\end{remark}


% ============================================================
\section{Cylinder Geometry of Division Words}
\label{sec:cylinder-geometry}
% ============================================================

Finite admissible division words determine exact \(2\)-adic realization
classes. For each admissible division word \(\omega\), exact realizability
determines a unique residue class modulo \(2^{D(\omega)+1}\), with a unique
canonical representative
\[
0\le r_\omega<2^{D(\omega)+1}.
\]

These residue classes naturally belong to the more general arithmetic
geometry of congruence classes modulo powers of \(2\). We first develop that
dyadic structure independently of the Collatz dynamics and then specialize
it to the realization classes associated with division words.

% ------------------------------------------------------------
\subsection{Dyadic Cylinders and Division-Word Cylinders}
% ------------------------------------------------------------

\begin{definition}[Dyadic Cylinder]
For \(r,k\in\mathbb N\), define the dyadic cylinder
\[
C(r,k)
:=
\left\{
x\in\mathbb N :
x\equiv r\pmod{2^k}
\right\}.
\]
\end{definition}

\begin{remark}
The residue \(r\) need not be canonical in this definition. When uniqueness
of the representative is required, we shall impose
\[
0\le r<2^k.
\]
The definition itself is purely arithmetic and is independent of the
Collatz map.
\end{remark}

\begin{definition}[Cylinder of a Division Word]
Let \(\omega\) be an admissible division word, and let \(r_\omega\) be its
canonical exact realization residue. Define
\[
C_\omega
:=
C\bigl(r_\omega,D(\omega)+1\bigr).
\]
Equivalently,
\[
C_\omega
=
\left\{
x\in\mathbb N :
x\equiv r_\omega
\pmod{2^{D(\omega)+1}}
\right\}.
\]
By Theorem~\ref{thm:exact-realization-residue-class},
\[
C_\omega
=
R_\omega
=
\{x\in\mathbb N:x\text{ realizes }\omega\}.
\]
\end{definition}

% ------------------------------------------------------------
\subsection{Geometry of Dyadic Cylinders}
% ------------------------------------------------------------

The basic geometry of dyadic cylinders is independent of the Collatz map.
It follows solely from the divisibility relation among powers of \(2\).

\begin{proposition}[Dyadic Cylinder Containment]
\label{prop:dyadic-cylinder-containment}

Let \(r,s,k,\ell\in\mathbb N\) with \(k\le\ell\). Then
\[
C(s,\ell)\subseteq C(r,k)
\]
if and only if
\[
s\equiv r\pmod{2^k}.
\]
\end{proposition}

\begin{proof}
Suppose first that
\[
C(s,\ell)\subseteq C(r,k).
\]
Since \(s\in C(s,\ell)\), the inclusion implies
\[
s\in C(r,k),
\]
and therefore
\[
s\equiv r\pmod{2^k}.
\]

Conversely, suppose
\[
s\equiv r\pmod{2^k},
\]
and let \(x\in C(s,\ell)\). Then
\[
x\equiv s\pmod{2^\ell}.
\]
Since \(k\le\ell\),
\[
2^k\mid2^\ell,
\]
so reduction modulo \(2^k\) gives
\[
x\equiv s\pmod{2^k}.
\]
Combining this with
\[
s\equiv r\pmod{2^k}
\]
yields
\[
x\equiv r\pmod{2^k}.
\]
Hence \(x\in C(r,k)\), proving
\[
C(s,\ell)\subseteq C(r,k).
\]
\end{proof}

\begin{corollary}[Nested-or-Disjoint Dyadic Cylinders]
\label{cor:dyadic-cylinders-nested-disjoint}

Any two dyadic cylinders are either nested or disjoint. More precisely,
for arbitrary \(r,s,k,\ell\in\mathbb N\), at least one of
\[
C(s,\ell)\subseteq C(r,k),
\qquad
C(r,k)\subseteq C(s,\ell),
\qquad
C(r,k)\cap C(s,\ell)=\varnothing
\]
holds.
\end{corollary}

\begin{proof}
Assume without loss of generality that \(k\le\ell\). If the two cylinders
are disjoint, there is nothing to prove.

Otherwise choose
\[
x\in C(r,k)\cap C(s,\ell).
\]
Then
\[
x\equiv r\pmod{2^k}
\]
and, after reducing the second congruence modulo \(2^k\),
\[
x\equiv s\pmod{2^k}.
\]
Hence
\[
s\equiv r\pmod{2^k}.
\]
Proposition~\ref{prop:dyadic-cylinder-containment} therefore gives
\[
C(s,\ell)\subseteq C(r,k).
\]
The case \(\ell\le k\) is symmetric.
\end{proof}

\begin{theorem}[Binary Refinement of a Dyadic Cylinder]
\label{thm:dyadic-cylinder-binary-refinement}

Let
\[
0\le r<2^k.
\]
Then
\[
C(r,k)
=
C(r,k+1)
\sqcup
C(r+2^k,k+1),
\]
where \(\sqcup\) denotes disjoint union.
\end{theorem}

\begin{proof}
A residue modulo \(2^{k+1}\) reducing to \(r\) modulo \(2^k\) has exactly
one of the two forms
\[
r
\qquad\text{or}\qquad
r+2^k.
\]
Therefore every element of \(C(r,k)\) lies in exactly one of
\[
C(r,k+1)
\qquad\text{and}\qquad
C(r+2^k,k+1).
\]

Conversely, both of these residues are congruent to \(r\) modulo \(2^k\),
so both child cylinders are contained in \(C(r,k)\).

The two children are disjoint because
\[
r\not\equiv r+2^k\pmod{2^{k+1}}.
\]
Hence their union is disjoint and equals the parent cylinder.
\end{proof}

\begin{remark}
Thus the complete family of dyadic cylinders forms a binary refinement
geometry. Passing from depth \(k\) to depth \(k+1\) resolves one additional
binary residue bit, and any two cylinders are either nested or disjoint.

The realization cylinders introduced above inherit this arithmetic geometry
after the specialization
\[
r=r_\omega,
\qquad
k=D(\omega)+1.
\]
\end{remark}

% ------------------------------------------------------------
\subsection{Counting and Natural Density}
% ------------------------------------------------------------

For a canonical dyadic cylinder \(C(r,k)\), define
\[
A_{r,k}(N)
:=
\#\{x\in\mathbb N:0\le x<N,\ x\in C(r,k)\}.
\]

With the convention \(0\in\mathbb N\), this definition also treats the
zero residue class uniformly. In particular, the initial block
\[
[0,2^k)
\]
contains exactly one representative of every residue class modulo \(2^k\),
including \(0\). Cylinder membership is periodic with period \(2^k\).
Moreover, if
\[
0\le r<2^k,
\]
then every complete block of \(2^k\) consecutive natural numbers contains
exactly one member of \(C(r,k)\).

\begin{theorem}[Dyadic Cylinder Counting Bounds]
\label{thm:dyadic-cylinder-counting}

Let
\[
0\le r<2^k.
\]
Then for every \(N\in\mathbb N\),
\[
\left\lfloor\frac{N}{2^k}\right\rfloor
\le
A_{r,k}(N)
\le
\left\lfloor\frac{N}{2^k}\right\rfloor+1.
\]
\end{theorem}

\begin{proof}
Write
\[
N=q\,2^k+s,
\qquad
0\le s<2^k.
\]
The first \(q\,2^k\) natural numbers consist of \(q\) complete periods of
length \(2^k\), each containing exactly one member of the cylinder.
They therefore contribute exactly \(q\) points.

The remaining interval has length \(s<2^k\), so it contains at most one
additional member of the same residue class. Hence
\[
q\le A_{r,k}(N)\le q+1.
\]
Since
\[
q=\left\lfloor\frac{N}{2^k}\right\rfloor,
\]
the result follows.
\end{proof}

\begin{theorem}[Natural Density of a Dyadic Cylinder]
\label{thm:dyadic-cylinder-density}

Every canonical dyadic cylinder \(C(r,k)\) has ordinary natural density
\[
2^{-k}.
\]
More precisely, for \(N>0\),
\[
\left|
\frac{A_{r,k}(N)}{N}
-
\frac{1}{2^k}
\right|
\le
\frac{1}{N},
\]
and therefore
\[
\lim_{N\to\infty}
\frac{A_{r,k}(N)}{N}
=
\frac{1}{2^k}.
\]
\end{theorem}

\begin{proof}
Write
\[
N=q\,2^k+s,
\qquad
0\le s<2^k.
\]
By Theorem~\ref{thm:dyadic-cylinder-counting},
\[
A_{r,k}(N)\in\{q,q+1\}.
\]
In either case,
\[
\left|
A_{r,k}(N)-\frac{N}{2^k}
\right|
\le1.
\]
Dividing by \(N>0\) gives
\[
\left|
\frac{A_{r,k}(N)}{N}
-
\frac{1}{2^k}
\right|
\le
\frac1N.
\]
Since \(1/N\to0\), the stated density follows.
\end{proof}

\begin{corollary}[Density of a Realization Cylinder]
\label{cor:realization-cylinder-density}

Let \(\omega\) be an admissible division word. Then
\[
C_\omega
=
C\bigl(r_\omega,D(\omega)+1\bigr)
\]
has ordinary natural density
\[
\boxed{
\operatorname{dens}(C_\omega)
=
2^{-(D(\omega)+1)}.
}
\]
\end{corollary}

\begin{proof}
The canonical realization residue satisfies
\[
0\le r_\omega<2^{D(\omega)+1}.
\]
Apply Theorem~\ref{thm:dyadic-cylinder-density} with
\[
r=r_\omega,
\qquad
k=D(\omega)+1.
\]
\end{proof}

% ------------------------------------------------------------
\subsection{Realization Cylinder Forest}
% ------------------------------------------------------------

We now specialize the generic dyadic geometry to the cylinders determined
by admissible division words.

Write
\[
\omega\preceq\eta
\]
when \(\omega\) is a prefix of \(\eta\).

\begin{lemma}[Proper Prefixes Increase Dyadic Depth]
\label{lem:proper-prefix-division-count}

Let \(\omega\) be a proper prefix of an admissible division word \(\eta\).
Then
\[
D(\omega)<D(\eta).
\]
\end{lemma}

\begin{proof}
Since \(\omega\) is a proper prefix of \(\eta\), write
\[
\eta=\omega\xi
\]
with \(\xi\) nonempty. Every entry of \(\eta\), and hence every entry of
\(\xi\), is positive. Therefore
\[
D(\xi)>0.
\]
By additivity of total division count under concatenation,
\[
D(\eta)=D(\omega)+D(\xi),
\]
and hence
\[
D(\omega)<D(\eta).
\]
\end{proof}

\begin{theorem}[Exact Cylinder Inclusion Criterion]
\label{thm:cylinder-inclusion}

Let \(\omega\) and \(\eta\) be admissible division words. Then
\[
C_\eta\subseteq C_\omega
\]
if and only if
\[
\omega\preceq\eta.
\]
Equivalently,
\[
\boxed{
\omega\preceq\eta
\quad\Longleftrightarrow\quad
C_\eta\subseteq C_\omega.
}
\]
\end{theorem}

\begin{proof}
Suppose first that
\[
C_\eta\subseteq C_\omega.
\]
Since
\[
r_\eta\in C_\eta,
\]
the assumed inclusion gives
\[
r_\eta\in C_\omega.
\]
Thus \(r_\eta\) realizes both \(\eta\) and \(\omega\). By
Lemma~\ref{lem:common-realization-prefix}, the two words are
prefix-comparable.

If \(\omega\preceq\eta\), we are done. Suppose instead that \(\eta\) is a
proper prefix of \(\omega\). By
Lemma~\ref{lem:proper-prefix-division-count},
\[
D(\eta)<D(\omega).
\]

Now
\[
r_\eta+2^{D(\eta)+1}
\]
belongs to the same cylinder \(C_\eta\) as \(r_\eta\). By the assumed
inclusion, both numbers therefore belong to \(C_\omega\). Since
\(C_\omega\) is a single residue class modulo \(2^{D(\omega)+1}\), the two
numbers are congruent to one another modulo \(2^{D(\omega)+1}\). Hence
\[
2^{D(\omega)+1}
\mid
2^{D(\eta)+1}.
\]
But
\[
D(\eta)<D(\omega),
\]
and therefore
\[
0<2^{D(\eta)+1}<2^{D(\omega)+1},
\]
making this divisibility impossible. Thus the reverse prefix relation
cannot be proper, and consequently
\[
\omega\preceq\eta.
\]

Conversely, suppose
\[
\omega\preceq\eta.
\]
Every natural number realizing \(\eta\) also realizes each prefix of
\(\eta\), and hence realizes \(\omega\). Since the cylinders are exactly
the corresponding realization sets,
\[
C_\eta\subseteq C_\omega.
\]
\end{proof}

\begin{corollary}[Forest Consequences]
\label{cor:cylinder-forest-consequences}

Let \(\omega\) and \(\eta\) be admissible division words.

\begin{enumerate}
\item
\[
C_\omega=C_\eta
\quad\Longleftrightarrow\quad
\omega=\eta.
\]

\item
\[
C_\eta\subsetneq C_\omega
\quad\Longleftrightarrow\quad
\omega\text{ is a proper prefix of }\eta.
\]

\item If neither word is a prefix of the other, then
\[
C_\omega\cap C_\eta=\varnothing.
\]

\item Distinct admissible words of equal length have disjoint cylinders.
\end{enumerate}
\end{corollary}

\begin{proof}
The first two statements follow immediately from
Theorem~\ref{thm:cylinder-inclusion}.

For the third, generic dyadic cylinders are nested or disjoint. If
\(C_\omega\) and \(C_\eta\) were not disjoint, one would be contained in
the other, and Theorem~\ref{thm:cylinder-inclusion} would force one word
to be a prefix of the other.

The fourth statement follows because two distinct words of equal length
cannot be prefix-comparable.
\end{proof}

\begin{remark}[Forest Geometry]
Theorem~\ref{thm:cylinder-inclusion} identifies symbolic ancestry exactly
with arithmetic containment, with the order reversed:
\[
\omega\preceq\eta
\quad\Longleftrightarrow\quad
C_\eta\subseteq C_\omega.
\]

Since admissible division words are nonempty, there is no empty-word root.
The length-one admissible words therefore form the roots of a family of
disjoint prefix trees, giving a natural realization-cylinder forest.

An edge of the symbolic forest appends one additional division count.
Because that division count may exceed \(1\), such an edge need not
correspond to a single binary dyadic refinement step: it may descend
several dyadic levels at once. Nevertheless every proper symbolic
descendant determines a strict arithmetic refinement to a cylinder of
strictly greater dyadic depth.

Thus the realization forest is embedded in the universal binary geometry
of dyadic cylinders, with its symbolic prefix order represented exactly
by cylinder containment.
\end{remark}


% ============================================================
\section{Affine Contraction and Descent Cylinders}
\label{sec:affine-contraction-descent}
% ============================================================

The preceding sections separated three distinct structures associated with
a division word: its formal affine map, its exact realization class, and
the dyadic cylinder geometry of that class. We now connect the affine
slope to actual descent along realized trajectories.

Recall that
\[
F_\omega(x)
=
\frac{3^{|\omega|}x+c(\omega)}
     {2^{D(\omega)}}
\]
has affine slope
\[
\lambda_\omega
=
\frac{3^{|\omega|}}{2^{D(\omega)}}.
\]
The dyadic--ternary threshold established earlier gives
\[
\lambda_\omega<1
\quad\Longleftrightarrow\quad
D(\omega)\ge B(|\omega|).
\]

A subunit slope alone does not imply
\[
F_\omega(x)<x,
\]
because the affine constant \(c(\omega)\) is positive. The key point is
that this obstruction is finite: once the initial value is sufficiently
large, the negative contribution from the subunit linear coefficient
dominates the fixed affine term.

% ------------------------------------------------------------
\subsection{Affine Contraction and Pointwise Descent}
% ------------------------------------------------------------

\begin{proposition}[Descent Threshold for Subunit Slope]
\label{prop:affine-descent-threshold}

Let \(\omega\) be a nonempty division word satisfying
\[
2^{D(\omega)}>3^{|\omega|}.
\]
Then for every \(x\in\mathbb Q\),
\[
F_\omega(x)<x
\]
if and only if
\[
\bigl(2^{D(\omega)}-3^{|\omega|}\bigr)x
>
c(\omega).
\]

Equivalently,
\[
F_\omega(x)<x
\quad\Longleftrightarrow\quad
x>
\frac{c(\omega)}
     {2^{D(\omega)}-3^{|\omega|}}.
\]
\end{proposition}

\begin{proof}
Starting from
\[
F_\omega(x)
=
\frac{3^{|\omega|}x+c(\omega)}
     {2^{D(\omega)}},
\]
and using \(2^{D(\omega)}>0\), we have
\[
F_\omega(x)<x
\]
if and only if
\[
3^{|\omega|}x+c(\omega)
<
2^{D(\omega)}x.
\]
Rearranging gives
\[
c(\omega)
<
\bigl(2^{D(\omega)}-3^{|\omega|}\bigr)x.
\]

Since
\[
2^{D(\omega)}-3^{|\omega|}>0,
\]
division by this positive quantity yields the equivalent threshold
\[
x>
\frac{c(\omega)}
     {2^{D(\omega)}-3^{|\omega|}}.
\]
\end{proof}

\begin{corollary}[Eventual Descent on a Realization Cylinder]
\label{cor:eventual-descent-on-cylinder}

Let \(\omega\) be an admissible division word with
\[
\lambda_\omega<1.
\]
Then all sufficiently large \(x\in C_\omega\) satisfy
\[
f^{|\omega|}(x)<x.
\]
Consequently, only finitely many members of \(C_\omega\) can fail to
lie below their starting value at the endpoint of the trajectory segment
encoded by \(\omega\).
\end{corollary}

\begin{proof}
If \(x\in C_\omega\), then \(x\) realizes the admissible word
\(\omega\). Hence \(x\) is odd and positive, so the iterates
\(f^i(x)\) are genuine fully accelerated Collatz iterates. Moreover,
\[
F_\omega(x)=f^{|\omega|}(x).
\]
Since \(\lambda_\omega<1\),
\[
2^{D(\omega)}>3^{|\omega|}.
\]
By Proposition~\ref{prop:affine-descent-threshold},
\[
F_\omega(x)<x
\]
for every
\[
x>
\frac{c(\omega)}
     {2^{D(\omega)}-3^{|\omega|}}.
\]
Only finitely many natural numbers lie below this fixed bound, so only
finitely many members of \(C_\omega\) can fail the stated descent
inequality.
\end{proof}

\begin{remark}
This is the precise distinction between affine slope contraction and
pointwise descent.

The condition
\[
\lambda_\omega<1
\]
depends only on the multiplicative balance
\[
3^{|\omega|}
\quad\text{versus}\quad
2^{D(\omega)}.
\]
Pointwise descent additionally depends on the affine constant
\(c(\omega)\). But for a fixed subunit-slope word, that constant contributes
only the finite threshold
\[
\frac{c(\omega)}
     {2^{D(\omega)}-3^{|\omega|}}.
\]

Thus the affine term can delay descent for finitely many small realizations,
but it cannot alter the asymptotic behavior of the realization cylinder.
The next subsection makes this density statement precise.
\end{remark}

% ------------------------------------------------------------
\subsection{Asymptotic Descent Cylinders}
% ------------------------------------------------------------

Let \(\omega\) be an admissible division word with
\[
\lambda_\omega<1.
\]
Define its descending realization set by
\[
C_\omega^\downarrow
:=
\left\{
x\in C_\omega :
f^{|\omega|}(x)<x
\right\}.
\]
Since every \(x\in C_\omega\) realizes \(\omega\),
\[
f^{|\omega|}(x)=F_\omega(x),
\]
and hence equivalently
\[
C_\omega^\downarrow
=
\left\{
x\in C_\omega :
F_\omega(x)<x
\right\}.
\]

Let
\[
E_\omega
:=
C_\omega\setminus C_\omega^\downarrow
\]
denote the corresponding set of non-descending realizations.

\begin{lemma}[Finite Affine Exceptional Set]
\label{lem:finite-affine-exceptional-set}

Let \(\omega\) be admissible and satisfy
\[
\lambda_\omega<1.
\]
Then every \(x\in E_\omega\) satisfies
\[
\bigl(2^{D(\omega)}-3^{|\omega|}\bigr)x
\le
c(\omega).
\]
In particular,
\[
x\le c(\omega),
\]
so \(E_\omega\) is finite.
\end{lemma}

\begin{proof}
If \(x\in E_\omega\), then
\[
F_\omega(x)\ge x.
\]
By Proposition~\ref{prop:affine-descent-threshold}, this is equivalent to
\[
\bigl(2^{D(\omega)}-3^{|\omega|}\bigr)x
\le
c(\omega).
\]

Since \(\lambda_\omega<1\),
\[
2^{D(\omega)}-3^{|\omega|}>0.
\]
The difference is a positive integer, and therefore at least \(1\). Hence
\[
x
\le
\bigl(2^{D(\omega)}-3^{|\omega|}\bigr)x
\le
c(\omega).
\]
Thus
\[
E_\omega
\subseteq
\{0,1,\dots,c(\omega)\},
\]
and \(E_\omega\) is finite.
\end{proof}

\begin{theorem}[Density Preservation under Descent]
\label{thm:density-preservation-eventual-descent}

Let \(\omega\) be an admissible division word satisfying
\[
\lambda_\omega<1.
\]
Then
\[
\operatorname{dens}(C_\omega^\downarrow)
=
\operatorname{dens}(C_\omega)
=
2^{-(D(\omega)+1)}.
\]
\end{theorem}

\begin{proof}
By Corollary~\ref{cor:realization-cylinder-density},
\[
\operatorname{dens}(C_\omega)
=
2^{-(D(\omega)+1)}.
\]

By Lemma~\ref{lem:finite-affine-exceptional-set}, the exceptional set
\(E_\omega\) is finite. Hence
\[
\frac{\#(E_\omega\cap[0,N))}{N}
\longrightarrow0.
\]

Since
\[
C_\omega
=
C_\omega^\downarrow
\sqcup
E_\omega,
\]
for every \(N\),
\[
\#(C_\omega\cap[0,N))
=
\#(C_\omega^\downarrow\cap[0,N))
+
\#(E_\omega\cap[0,N)).
\]
Dividing by \(N\) and passing to the limit gives
\[
\operatorname{dens}(C_\omega^\downarrow)
=
\operatorname{dens}(C_\omega)
=
2^{-(D(\omega)+1)}.
\]
\end{proof}

\begin{remark}
This obstruction is asymptotically negligible: the non-descending
exceptional set \(E_\omega\) is finite, so removing it does not change
natural density:
\[
\operatorname{dens}(C_\omega^\downarrow)
=
\operatorname{dens}(C_\omega).
\]

Thus every admissible subunit-slope cylinder contributes its full dyadic
density
\[
2^{-(D(\omega)+1)}
\]
to accelerated descent, up to a finite exceptional set.
\end{remark}

\begin{remark}[From Local Descent to Global Residual Structure]
The preceding theorem is a statement about a fixed admissible division word.
Whenever
\[
\lambda_\omega<1,
\]
the corresponding realization cylinder contributes its full natural density
to accelerated descent, up to a finite exceptional set.

This is not yet a global density statement. There are infinitely many
admissible division words at every positive symbolic length, so the argument
cannot proceed by counting all admissible cylinders and comparing their
number with the density of a typical cylinder.

Instead, the global argument must keep track of those trajectories whose
prefixes have not yet entered the subunit-slope regime. The next section
introduces the resulting first-contraction and residual cylinder structure.
Its purpose is to partition the cylinder forest according to the first prefix
at which affine contraction occurs and to estimate the dyadic mass that
remains unresolved after successive refinements.
\end{remark}


% ============================================================
\section{First-Contraction and the Residual Forest}
\label{sec:residual-forest}
% ============================================================

The preceding section showed that an admissible division word with subunit
affine slope contributes its full realization-cylinder density to
accelerated descent, apart from finitely many exceptional initial values.
To obtain a global statement, it remains to organize trajectories according
to the first prefix at which the affine slope becomes strictly less than \(1\).

This requires a restricted symbolic structure. The full family of admissible
division words is infinite at every positive word length and therefore does
not provide a finite horizontal decomposition. Instead, we retain only those
prefixes whose total division count has not yet reached the dyadic threshold
\(B(j)\). These prefixes form the \emph{residual forest}.

For a division word \(\omega\) and \(1\le j\le|\omega|\), write
\[
\omega_{\le j}
\]
for its prefix of length \(j\). Recall that
\[
B(j)
\]
is the least integer \(k\) satisfying
\[
2^k>3^j.
\]
Thus
\[
D(\omega_{\le j})\ge B(j)
\]
is equivalent to the affine slope of the prefix \(\omega_{\le j}\) being
strictly less than \(1\).

% ------------------------------------------------------------
\subsection{First-Contraction and Residual Words}
% ------------------------------------------------------------

\begin{definition}[Residual Word]
\label{def:residual-word}

An admissible division word \(\omega\) is called \emph{residual} if every
nonempty prefix has total division count strictly below the dyadic threshold:
\[
D(\omega_{\le j})<B(j),
\qquad
1\le j\le|\omega|.
\]
Equivalently, every nonempty prefix has affine slope strictly greater than
\(1\).
\end{definition}

\begin{definition}[First-Contraction Word]
\label{def:first-contraction-word}

An admissible division word \(\omega\) of length \(m\) is called a
\emph{first-contraction word} if
\[
D(\omega)\ge B(m),
\]
while every proper nonempty prefix remains residual:
\[
D(\omega_{\le j})<B(j),
\qquad
1\le j<m.
\]
Thus the total division count of \(\omega\) reaches the dyadic threshold for
the first time at its terminal symbolic position.
\end{definition}

\begin{lemma}[Residual Prefix Closure]
\label{lem:residual-prefix-closure}

Every nonempty prefix of a residual word is residual.
\end{lemma}

\begin{proof}
Let \(\omega\) be residual and let
\[
\eta=\omega_{\le j},
\qquad
1\le j\le|\omega|.
\]
Admissibility is inherited by nonempty prefixes. Moreover, for every
\(1\le k\le j\),
\[
\eta_{\le k}=\omega_{\le k}.
\]
Since \(\omega\) is residual,
\[
D(\omega_{\le k})<B(k).
\]
Hence
\[
D(\eta_{\le k})<B(k)
\]
for every nonempty prefix of \(\eta\), so \(\eta\) is residual.
\end{proof}

\begin{proposition}[First-Contraction Antichain]
\label{prop:first-contraction-antichain}

Distinct first-contraction words are prefix-incomparable. Consequently,
their realization cylinders are pairwise disjoint, even when the words have
different symbolic lengths.
\end{proposition}

\begin{proof}
Suppose that distinct first-contraction words \(\omega\) and \(\eta\)
satisfied
\[
\omega\prec\eta,
\]
where \(\prec\) denotes the proper-prefix relation.

Since \(\omega\) is a first-contraction word,
\[
D(\omega)\ge B(|\omega|).
\]
But \(\omega\) is a proper prefix of the first-contraction word \(\eta\), so
the defining residual-prefix condition for \(\eta\) requires
\[
D(\omega)<B(|\omega|),
\]
a contradiction.

Thus neither word can be a proper prefix of the other. By the
symbolic--arithmetic correspondence for realization cylinders,
prefix-incomparable division words have disjoint cylinders.
\end{proof}

\begin{remark}
The antichain property is important, but it does not by itself license an
infinite sum of cylinder densities. The family of first-contraction words
at a fixed symbolic length is generally infinite, since the total division
count at the terminal depth may exceed the dyadic threshold at that length
by an arbitrarily large amount.

The finite objects used below are instead the \emph{residual levels}.
\end{remark}

% ------------------------------------------------------------
\subsection{Residual Levels and Exit Layers}
% ------------------------------------------------------------

For \(m\ge1\), define the residual word level
\[
\mathcal W_m^{\mathrm{res}}
:=
\left\{
\omega:
|\omega|=m
\text{ and }
\omega\text{ is residual}
\right\}.
\]

\begin{proposition}[Finiteness of Residual Levels]
\label{prop:residual-level-finite}

For every \(m\ge1\), the set
\[
\mathcal W_m^{\mathrm{res}}
\]
is finite.
\end{proposition}

\begin{proof}
If
\[
\omega=(d_0,\dots,d_{m-1})
\in\mathcal W_m^{\mathrm{res}},
\]
then residuality at the terminal prefix gives
\[
D(\omega)<B(m).
\]
Since every \(d_i\) is positive and
\[
d_i\le D(\omega),
\]
we obtain
\[
1\le d_i<B(m)
\]
for every \(i\).

Thus every residual word of length \(m\) is a length-\(m\) word over the
finite alphabet
\[
\{1,2,\dots,B(m)-1\}.
\]
Only finitely many such words exist.
\end{proof}

\begin{corollary}[Root of the Residual Forest]
\label{cor:residual-root}

The first residual level consists of the single word
\[
\mathcal W_1^{\mathrm{res}}=\{(1)\}.
\]
\end{corollary}

\begin{proof}
A residual word of length one has the form \((d)\) with \(d\ge1\), while
residuality requires
\[
d<B(1)=2.
\]
Hence \(d=1\).
\end{proof}

The first-contraction words of length one are therefore precisely the
one-symbol words \((d)\) with \(d\ge2\). Their realizations lie outside the
initial residual set \(\mathcal R_1\), so the residual-forest decomposition
below begins with the exit layer \(\mathcal F_2\). The complementary
length-one population is handled separately in the final descent argument.

For \(m\ge1\), define the residual realization set
\[
\mathcal R_m
:=
\bigcup_{\omega\in\mathcal W_m^{\mathrm{res}}}
C_\omega.
\]
Thus \(\mathcal R_m\) consists precisely of starting values whose first
\(m\) division counts remain residual.

\begin{theorem}[Nested Residual Forest]
\label{thm:residual-levels-nested}

For every \(m\ge1\),
\[
\mathcal R_{m+1}\subseteq\mathcal R_m.
\]
\end{theorem}

\begin{proof}
Let \(x\in\mathcal R_{m+1}\). Then \(x\) realizes some residual word
\(\omega\) of length \(m+1\).

By Lemma~\ref{lem:residual-prefix-closure}, the prefix
\[
\omega_{\le m}
\]
is residual. Since realization of a word implies realization of each of its
prefixes, \(x\) realizes \(\omega_{\le m}\). Therefore
\[
x\in\mathcal R_m.
\]
\end{proof}

The set of starting values whose realized trajectories leave the residual
forest at the next refinement step is now defined by
\[
\mathcal F_{m+1}
:=
\mathcal R_m\setminus\mathcal R_{m+1},
\qquad
m\ge1.
\]

\begin{theorem}[Exit Layer Equals First Contraction]
\label{thm:exit-layer-first-contraction}

For every \(m\ge1\),
\[
\mathcal F_{m+1}
=
\left\{
x\in\mathbb N:
x\text{ realizes a first-contraction word of length }m+1
\right\}.
\]
\end{theorem}

\begin{proof}
Suppose first that
\[
x\in\mathcal F_{m+1}.
\]
Then
\[
x\in\mathcal R_m
\qquad\text{but}\qquad
x\notin\mathcal R_{m+1}.
\]
Hence \(x\) realizes a residual word \(\omega\) of length \(m\).
In particular, \(\omega\) is admissible, so \(x\) is odd and positive.

Let \(\eta\) denote the actual division word of length \(m+1\) generated by
the trajectory of \(x\). By
Lemma~\ref{lem:generated-division-words}, \(\eta\) is admissible. Its
length-\(m\) prefix must equal \(\omega\), because a fixed starting value
determines a unique realized word of each length.

Therefore every proper nonempty prefix of \(\eta\) is residual. If the full
word \(\eta\) also satisfied
\[
D(\eta)<B(m+1),
\]
then \(\eta\) itself would be residual, and \(x\) would belong to
\(\mathcal R_{m+1}\), contrary to assumption. Hence
\[
D(\eta)\ge B(m+1),
\]
so \(\eta\) is a first-contraction word.

Conversely, suppose that \(x\) realizes a first-contraction word \(\eta\) of
length \(m+1\). Its length-\(m\) prefix is residual, so
\[
x\in\mathcal R_m.
\]

If \(x\in\mathcal R_{m+1}\), then \(x\) would also realize some residual word
\(\rho\) of length \(m+1\). But two words of the same length realized by the
same starting value must coincide, so
\[
\rho=\eta.
\]
This is impossible, because \(\rho\) is residual at its terminal depth while
\(\eta\) reaches the contraction threshold there. Thus
\[
x\notin\mathcal R_{m+1},
\]
and hence
\[
x\in\mathcal F_{m+1}.
\]
\end{proof}

\begin{corollary}[Finite-Stage Forest Decomposition]
\label{cor:residual-forest-decomposition}

For every \(m\ge1\),
\[
\boxed{
\mathcal R_m
=
\mathcal F_{m+1}
\sqcup
\mathcal R_{m+1}.
}
\]
\end{corollary}

\begin{proof}
By definition,
\[
\mathcal F_{m+1}
=
\mathcal R_m\setminus\mathcal R_{m+1},
\]
and Theorem~\ref{thm:residual-levels-nested} gives
\[
\mathcal R_{m+1}\subseteq\mathcal R_m.
\]
The two sets are therefore disjoint and their union is exactly
\(\mathcal R_m\).
\end{proof}

\begin{remark}
This decomposition is the structural conservation law of the residual
forest. Every trajectory that is residual through depth \(m\) has exactly
one of two fates at the next symbolic step:

\[
\text{first contraction at depth }m+1,
\]
or
\[
\text{continued residual survival through depth }m+1.
\]

No trajectory is lost and none is counted twice.
\end{remark}

\begin{definition}[Infinite Residual Set]
\label{def:infinite-residual-set}

Define
\[
\mathcal R_\infty
:=
\bigcap_{m\ge1}\mathcal R_m.
\]
Thus \(\mathcal R_\infty\) consists of those starting values whose realized
division-word prefixes remain residual at every finite depth.
\end{definition}

\begin{corollary}[Global First-Contraction Decomposition]
\label{cor:global-first-contraction-decomposition}

The initial residual set decomposes disjointly as
\[
\boxed{
\mathcal R_1
=
\left(
\bigsqcup_{m\ge1}\mathcal F_{m+1}
\right)
\sqcup
\mathcal R_\infty.
}
\]
\end{corollary}

\begin{proof}
Let \(x\in\mathcal R_1\). Either
\[
x\in\mathcal R_m
\]
for every \(m\ge1\), in which case
\[
x\in\mathcal R_\infty,
\]
or else there is a least \(m\ge1\) such that
\[
x\notin\mathcal R_{m+1}.
\]
By minimality,
\[
x\in\mathcal R_m,
\]
and hence
\[
x\in
\mathcal R_m\setminus\mathcal R_{m+1}
=
\mathcal F_{m+1}.
\]

Conversely, every exit layer \(\mathcal F_{m+1}\) lies in
\(\mathcal R_1\) by nestedness, as does \(\mathcal R_\infty\).

The exit layers are pairwise disjoint because an element has at most one
first level at which it leaves the residual forest, and no exit-layer
element belongs to \(\mathcal R_\infty\). Hence the decomposition is
disjoint.
\end{proof}

\begin{remark}[Symbolic Conservation]
Every trajectory in the initial residual set has exactly one of two
possibilities: it reaches a first-contraction layer at some finite symbolic
depth, or it remains residual at every finite depth.

Thus symbolic refinement neither creates nor duplicates trajectories. It
partitions the initial residual set into first-contraction exits at finite
symbolic depths and the infinite-depth residual set.
\end{remark}

% ------------------------------------------------------------
\subsection{Residual Mass and Finite-Stage Conservation}
% ------------------------------------------------------------

Because \(\mathcal W_m^{\mathrm{res}}\) is finite and distinct words at the
same symbolic length have disjoint realization cylinders, define the
residual mass at depth \(m\) by
\[
M_m
:=
\sum_{\omega\in\mathcal W_m^{\mathrm{res}}}
2^{-(D(\omega)+1)}.
\]

Each term is exactly the natural density of the corresponding realization
cylinder \(C_\omega\).

\begin{theorem}[Residual Mass Equals Residual Density]
\label{thm:residual-mass-density}

For every \(m\ge1\),
\[
\boxed{
\operatorname{dens}(\mathcal R_m)=M_m.
}
\]
\end{theorem}

\begin{proof}
For \(N\ge1\), let
\[
A_m(N)
:=
\#\bigl(\mathcal R_m\cap[0,N)\bigr).
\]
Since the cylinders
\[
\{C_\omega:\omega\in\mathcal W_m^{\mathrm{res}}\}
\]
form a finite pairwise-disjoint family,
\[
A_m(N)
=
\sum_{\omega\in\mathcal W_m^{\mathrm{res}}}
A_{r_\omega,D(\omega)+1}(N),
\]
where \(A_{r,k}(N)\) is the dyadic-cylinder counting function introduced
earlier.

For each residual word \(\omega\),
\[
\frac{A_{r_\omega,D(\omega)+1}(N)}{N}
\longrightarrow
2^{-(D(\omega)+1)}.
\]
Because the residual level is finite, the limits may be summed:
\[
\frac{A_m(N)}{N}
\longrightarrow
\sum_{\omega\in\mathcal W_m^{\mathrm{res}}}
2^{-(D(\omega)+1)}
=
M_m.
\]
Hence
\[
\operatorname{dens}(\mathcal R_m)=M_m.
\]
\end{proof}

\begin{theorem}[Exit-Layer Density]
\label{thm:first-contraction-layer-density}

For every \(m\ge1\), the first-contraction layer \(\mathcal F_{m+1}\) has
natural density
\[
\boxed{
\operatorname{dens}(\mathcal F_{m+1})
=
M_m-M_{m+1}.
}
\]
Consequently,
\[
M_{m+1}\le M_m.
\]
\end{theorem}

\begin{proof}
From Corollary~\ref{cor:residual-forest-decomposition},
\[
\mathcal R_m
=
\mathcal F_{m+1}
\sqcup
\mathcal R_{m+1}.
\]
Therefore, for every \(N\),
\[
\#(\mathcal R_m\cap[0,N))
=
\#(\mathcal F_{m+1}\cap[0,N))
+
\#(\mathcal R_{m+1}\cap[0,N)).
\]
Rearranging and dividing by \(N\) gives
\[
\frac{\#(\mathcal F_{m+1}\cap[0,N))}{N}
=
\frac{\#(\mathcal R_m\cap[0,N))}{N}
-
\frac{\#(\mathcal R_{m+1}\cap[0,N))}{N}.
\]
Passing to the limit and applying
Theorem~\ref{thm:residual-mass-density} yields
\[
\operatorname{dens}(\mathcal F_{m+1})
=
M_m-M_{m+1}.
\]
A natural density is nonnegative, so
\[
M_{m+1}\le M_m.
\]
\end{proof}

\begin{corollary}[Finite-Stage Mass Conservation]
\label{cor:finite-stage-mass-conservation}

For \(m\ge1\) and \(k\ge1\),
\[
\boxed{
\sum_{j=0}^{k-1}
\operatorname{dens}(\mathcal F_{m+j+1})
=
M_m-M_{m+k}.
}
\]
Equivalently,
\[
M_m
=
\sum_{j=0}^{k-1}
\operatorname{dens}(\mathcal F_{m+j+1})
+
M_{m+k}.
\]
\end{corollary}

\begin{proof}
By Theorem~\ref{thm:first-contraction-layer-density},
\[
\operatorname{dens}(\mathcal F_{m+j+1})
=
M_{m+j}-M_{m+j+1}.
\]
Summing over \(0\le j<k\) gives the telescoping identity
\[
\sum_{j=0}^{k-1}
\bigl(M_{m+j}-M_{m+j+1}\bigr)
=
M_m-M_{m+k}.
\]
\end{proof}

\begin{remark}[Finite-Stage Versus Asymptotic Collapse]
The preceding identities are entirely finite-stage statements. They show
that residual mass is redistributed exactly between successive
first-contraction exit layers and the portion that remains unresolved.

No assertion has yet been made that
\[
M_m\longrightarrow0.
\]
In particular, no continuity-from-above property of natural density and no
countable additivity over infinitely many first-contraction cylinders is
being invoked.

The remaining task is quantitative: prove directly that the residual mass
\(M_m\) tends to zero. That decay estimate is developed in the next section.
\end{remark}


% ============================================================
\section{Residual Mass Collapse}
\label{sec:residual-mass-collapse}
% ============================================================

The preceding section reduced the unresolved part of the symbolic dynamics
to the residual forest. At each depth \(m\), its natural density is the
finite residual mass
\[
M_m
=
\sum_{\omega\in\mathcal W_m^{\mathrm{res}}}
2^{-(D(\omega)+1)},
\]
and the finite-stage conservation law gives
\[
\operatorname{dens}(\mathcal F_{m+1})
=
M_m-M_{m+1}.
\]

The remaining question is therefore quantitative:

\[
\boxed{M_m\longrightarrow0\ ?}
\]

A direct estimate of \(M_m\) is inconvenient because extension of a word
changes its dyadic weight by a factor depending on the appended division
count. We therefore introduce an auxiliary weight for which all positive
extensions can be summed geometrically. The resulting one-step contraction
is then transferred back to the true dyadic residual mass.

% ------------------------------------------------------------
\subsection{Auxiliary Weighted Residual Mass}
% ------------------------------------------------------------

For a residual word \(\omega\), define its auxiliary weight by
\[
w(\omega)
:=
\left(\frac{3}{8}\right)^{D(\omega)}.
\]
The corresponding weighted residual mass at depth \(m\) is
\[
W_m
:=
\sum_{\omega\in\mathcal W_m^{\mathrm{res}}}
\left(\frac{3}{8}\right)^{D(\omega)}.
\]

The choice \(3/8\) is adapted to the extension structure of division words.
If a positive division count \(d\) is appended to \(\omega\), then
\[
D(\omega d)=D(\omega)+d,
\]
and hence
\[
w(\omega d)
=
w(\omega)
\left(\frac{3}{8}\right)^d.
\]
For every \(m\ge1\), all possible residual extensions lie within the
finite range \(1\le d<B(m+1)\), whose geometric weight satisfies
\[
\sum_{d=1}^{B(m+1)-1}
\left(\frac38\right)^d
<
\sum_{d=1}^{\infty}
\left(\frac38\right)^d
=
\frac{3/8}{1-3/8}
=
\frac35.
\]
Thus \(3/5\) is a uniform upper bound for the finite extension sum at every
residual depth.

\begin{theorem}[One-Step Weighted Contraction]
\label{thm:weighted-residual-contraction}

For every \(m\ge1\),
\[
\boxed{
W_{m+1}\le\frac35\,W_m.
}
\]
\end{theorem}

\begin{proof}
Every residual word \(\eta\) of length \(m+1\) has a unique decomposition
\[
\eta=\omega d,
\]
where \(\omega\) is its length-\(m\) prefix and \(d\ge1\) is its final
division count. By residual prefix closure, \(\omega\) is residual.

Because \(\eta\) is residual at depth \(m+1\),
\[
D(\eta)<B(m+1).
\]
Since
\[
D(\eta)=D(\omega)+d
\]
and \(D(\omega)\ge0\), it follows that
\[
1\le d<B(m+1).
\]

Thus the actual residual children of the words in
\(\mathcal W_m^{\mathrm{res}}\) form a subfamily of the finite collection
of extensions with
\[
1\le d<B(m+1).
\]
Since all weights are nonnegative,
\[
\begin{aligned}
W_{m+1}
&\le
\sum_{\omega\in\mathcal W_m^{\mathrm{res}}}
\sum_{d=1}^{B(m+1)-1}
\left(\frac38\right)^{D(\omega)+d}
\\[4pt]
&=
\sum_{\omega\in\mathcal W_m^{\mathrm{res}}}
\left(\frac38\right)^{D(\omega)}
\sum_{d=1}^{B(m+1)-1}
\left(\frac38\right)^d
\\[4pt]
&\le
\frac35
\sum_{\omega\in\mathcal W_m^{\mathrm{res}}}
\left(\frac38\right)^{D(\omega)}
\\[4pt]
&=
\frac35\,W_m.
\end{aligned}
\]
\end{proof}

At depth one, the only residual word is \((1)\). Therefore
\[
W_1=\frac38<\frac35.
\]

\begin{corollary}[Geometric Bound for Weighted Mass]
\label{cor:weighted-residual-geometric-bound}

For every \(m\ge1\),
\[
\boxed{
W_m\le\left(\frac35\right)^m.
}
\]
\end{corollary}

\begin{proof}
The result follows by induction from
Theorem~\ref{thm:weighted-residual-contraction} and the bound
\[
W_1\le\frac35.
\]
\end{proof}

\begin{remark}
The weighted mass \(W_m\) is auxiliary; it is not itself the natural density
of the residual level. Its purpose is to expose a uniformly subcritical
extension rule. The next step transfers this geometric decay to the genuine
dyadic mass \(M_m\).
\end{remark}

% ------------------------------------------------------------
\subsection{Transfer to the Dyadic Residual Mass}
% ------------------------------------------------------------

The exact density contribution of a realization cylinder with total
division count \(D\) is
\[
2^{-(D+1)}.
\]
This weight admits the factorization
\[
\boxed{
\frac{1}{2^{D+1}}
=
\frac12
\left(\frac43\right)^D
\left(\frac38\right)^D.
}
\]
Indeed,
\[
\frac43\cdot\frac38=\frac12.
\]

The final factor is precisely the auxiliary weight used in \(W_m\). The
remaining factor can be controlled uniformly across a residual level.

If
\[
\omega\in\mathcal W_m^{\mathrm{res}},
\]
then residuality at its terminal prefix gives
\[
D(\omega)<B(m),
\]
and therefore
\[
D(\omega)\le B(m)-1.
\]
Since \(4/3>1\),
\[
\left(\frac43\right)^{D(\omega)}
\le
\left(\frac43\right)^{B(m)-1}.
\]

\begin{proposition}[Weighted Envelope for Residual Mass]
\label{prop:residual-weighted-envelope}

For every \(m\ge1\),
\[
M_m
\le
\frac12
\left(\frac43\right)^{B(m)-1}
W_m.
\]
Consequently,
\[
\boxed{
0\le M_m\le E_m,
}
\]
where
\[
E_m
:=
\frac12
\left(\frac43\right)^{B(m)-1}
\left(\frac35\right)^m.
\]
\end{proposition}

\begin{proof}
For each
\(\omega\in\mathcal W_m^{\mathrm{res}}\),
\[
\begin{aligned}
2^{-(D(\omega)+1)}
&=
\frac12
\left(\frac43\right)^{D(\omega)}
\left(\frac38\right)^{D(\omega)}
\\
&\le
\frac12
\left(\frac43\right)^{B(m)-1}
\left(\frac38\right)^{D(\omega)}.
\end{aligned}
\]
Summing over the finite residual level gives
\[
M_m
\le
\frac12
\left(\frac43\right)^{B(m)-1}W_m.
\]
Applying Corollary~\ref{cor:weighted-residual-geometric-bound} yields
\[
M_m
\le
\frac12
\left(\frac43\right)^{B(m)-1}
\left(\frac35\right)^m
=
E_m.
\]
Nonnegativity of \(M_m\) is immediate from its definition.
\end{proof}

The envelope \(E_m\) contains two competing effects. The factor
\[
\left(\frac43\right)^{B(m)-1}
\]
grows with symbolic depth, while
\[
\left(\frac35\right)^m
\]
decays geometrically. It remains to show that the latter dominates.

% ------------------------------------------------------------
\subsection{Five-Step Envelope Contraction}
% ------------------------------------------------------------

The necessary control of the threshold \(B(m)\) follows from the elementary
integer inequality
\[
3^5=243<256=2^8.
\]

\begin{lemma}[Five-Step Threshold Growth]
\label{lem:five-step-threshold-growth}

For every \(m\ge0\),
\[
\boxed{
B(m+5)\le B(m)+8.
}
\]
\end{lemma}

\begin{proof}
By definition of \(B(m)\),
\[
3^m<2^{B(m)}.
\]
Multiplying this inequality by
\[
3^5<2^8
\]
gives
\[
3^{m+5}
<
2^{B(m)+8}.
\]
Since \(B(m+5)\) is the least exponent \(k\) for which
\[
3^{m+5}<2^k,
\]
we conclude that
\[
B(m+5)\le B(m)+8.
\]
\end{proof}

Define
\[
q
:=
\left(\frac43\right)^8
\left(\frac35\right)^5.
\]
A direct calculation gives
\[
q
=
\frac{65536}{84375}
<1.
\]

\begin{theorem}[Five-Step Envelope Contraction]
\label{thm:five-step-envelope-contraction}

For every \(m\ge1\),
\[
\boxed{
E_{m+5}\le qE_m.
}
\]
\end{theorem}

\begin{proof}
By Lemma~\ref{lem:five-step-threshold-growth},
\[
B(m+5)-1
\le
B(m)-1+8.
\]
Therefore
\[
\begin{aligned}
E_{m+5}
&=
\frac12
\left(\frac43\right)^{B(m+5)-1}
\left(\frac35\right)^{m+5}
\\
&\le
\frac12
\left(\frac43\right)^{B(m)-1+8}
\left(\frac35\right)^{m+5}
\\
&=
\left(\frac43\right)^8
\left(\frac35\right)^5
E_m
\\
&=
qE_m.
\end{aligned}
\]
\end{proof}

\begin{corollary}[Decay of the Residual-Mass Envelope]
\label{cor:residual-envelope-decay}

The envelope satisfies
\[
\boxed{
E_m\longrightarrow0
\qquad (m\to\infty).
}
\]
\end{corollary}

\begin{proof}
Every positive integer can be written uniquely as
\[
m=r+5k,
\qquad
r\in\{1,2,3,4,5\},
\qquad
k\ge0.
\]
Fix \(r\in\{1,2,3,4,5\}\). Iterating
Theorem~\ref{thm:five-step-envelope-contraction} gives
\[
E_{r+5k}
\le
q^k E_r.
\]
Since
\[
0<q<1,
\]
we have
\[
q^kE_r\longrightarrow0,
\]
and therefore
\[
E_{r+5k}\longrightarrow0
\]
for each \(r\in\{1,2,3,4,5\}\).

These five subsequences exhaust all positive indices. Hence
\[
E_m\longrightarrow0.
\]
\end{proof}

\begin{remark}
The five-step grouping is not an asymptotic approximation. It is an exact
integer comparison chosen so that the possible threshold growth and the
weighted-mass decay can be compared uniformly:
\[
3^5<2^8
\]
controls the former, while
\[
\left(\frac35\right)^5
\]
controls the latter. Their combined multiplier is strictly below \(1\).
\end{remark}

% ------------------------------------------------------------
\subsection{Residual Mass Collapse and First-Contraction Exhaustion}
% ------------------------------------------------------------

\begin{theorem}[Residual Mass Collapse]
\label{thm:residual-mass-collapse}

The residual mass tends to zero:
\[
\boxed{
M_m\longrightarrow0
\qquad (m\to\infty).
}
\]
Equivalently,
\[
\operatorname{dens}(\mathcal R_m)\longrightarrow0.
\]
\end{theorem}

\begin{proof}
For every \(m\ge1\),
Proposition~\ref{prop:residual-weighted-envelope} gives
\[
0\le M_m\le E_m.
\]
By Corollary~\ref{cor:residual-envelope-decay},
\[
E_m\longrightarrow0.
\]
The squeeze theorem therefore yields
\[
M_m\longrightarrow0.
\]
Since
\[
M_m=\operatorname{dens}(\mathcal R_m),
\]
the equivalent density statement follows.
\end{proof}

The finite-stage conservation law from the preceding section now has an
asymptotic consequence.

\begin{corollary}[Exhaustion by First-Contraction Mass]
\label{cor:first-contraction-mass-exhaustion}

For every \(k\ge1\),
\[
\sum_{j=0}^{k-1}
\operatorname{dens}(\mathcal F_{j+2})
=
M_1-M_{k+1},
\]
and therefore
\[
\boxed{
\lim_{k\to\infty}
\sum_{j=0}^{k-1}
\operatorname{dens}(\mathcal F_{j+2})
=
M_1.
}
\]
\end{corollary}

\begin{proof}
Finite-stage mass conservation gives
\[
\sum_{j=0}^{k-1}
\operatorname{dens}(\mathcal F_{j+2})
=
M_1-M_{k+1}.
\]
By Theorem~\ref{thm:residual-mass-collapse},
\[
M_{k+1}\longrightarrow0.
\]
Hence the partial sums converge to \(M_1\).
\end{proof}

\begin{remark}[Asymptotic Conservation]
The residual forest therefore loses no asymptotic mass except through its
first-contraction exit layers. At every finite stage,
\[
M_1
=
\sum_{j=0}^{k-1}
\operatorname{dens}(\mathcal F_{j+2})
+
M_{k+1},
\]
and the terminal residual term tends to zero.

Thus the finite partial sums of first-contraction layer densities
asymptotically exhaust the initial residual mass \(M_1\). This is a statement
about the limit of finite-stage density identities; it does not appeal to
countable additivity of natural density.
\end{remark}

\begin{remark}
The conclusion established here concerns affine first contraction. On each
fixed subunit-slope realization cylinder, the preceding section showed that
the exceptional set of non-descending realizations is finite.

This fixed-cylinder statement does not yet control an entire
first-contraction layer, which may contain infinitely many terminal words.
The remaining task is therefore to obtain a uniform descent cutoff at each
exit depth, combine that control with residual mass collapse, and finally
transfer accelerated descent to eventual descent under the ordinary Collatz
map.
\end{remark}


% ============================================================
\section{Density-One Eventual Descent}
\label{sec:density-one-descent}
% ============================================================

The preceding section proved that the residual mass tends to zero:
\[
M_m=\operatorname{dens}(\mathcal R_m)\longrightarrow0.
\]
It remains to convert this symbolic mass collapse into genuine dynamical
descent and then transfer the result from the fully accelerated map to the
ordinary Collatz map.

Two issues must be addressed. First, a first-contraction cylinder may contain
finitely many small realizations that do not yet descend below their starting
values. Since a fixed first-contraction layer may contain infinitely many
child words, these exceptional sets cannot simply be united word by word.
Instead, a uniform cutoff for each entire layer is required.

Second, the residual forest describes the fully accelerated dynamics of odd
positive integers. The final result concerns the ordinary Collatz map on
\(\mathbb N_{>0}\). The last part of the argument identifies the initial
residual class, handles the complementary congruence classes directly, and
transfers accelerated descent to ordinary Collatz descent.

% ------------------------------------------------------------
\subsection{Uniform Descent on First-Contraction Layers}
% ------------------------------------------------------------

Recall that
\[
\mathcal F_{m+1}
=
\mathcal R_m\setminus\mathcal R_{m+1}
\]
is precisely the set of natural numbers realizing a first-contraction word of
length \(m+1\).

Let
\[
\eta=\omega d
\]
be such a word, where \(\omega\) is its residual prefix of length \(m\) and
\(d\ge1\) is its final division count. The affine concatenation formula gives
\[
c(\eta)
=
3c(\omega)+2^{D(\omega)}c((d)).
\]
Since
\[
c((d))=1,
\]
we obtain
\[
\boxed{
c(\omega d)
=
3c(\omega)+2^{D(\omega)}.
}
\]
In particular, the affine constant of a one-symbol extension depends on its
residual parent but not on the value of the appended division count.
This observation supplies a uniform bound across an entire
first-contraction layer.

\begin{theorem}[Uniform Descent on a First-Contraction Layer]
\label{thm:uniform-layer-descent}

For every \(m\ge1\), there exists \(K_m\in\mathbb N\) such that
\[
x\in\mathcal F_{m+1},
\qquad
x\ge K_m
\]
implies
\[
f^{m+1}(x)<x.
\]
\end{theorem}

\begin{proof}
The residual parent level
\[
\mathcal W_m^{\mathrm{res}}
\]
is finite. If it is empty, then
\[
\mathcal F_{m+1}=\varnothing,
\]
and the result is immediate.

Otherwise define
\[
K_m
:=
1+
\max_{\omega\in\mathcal W_m^{\mathrm{res}}}
\left(
3c(\omega)+2^{D(\omega)}
\right).
\]
This maximum exists because
\(\mathcal W_m^{\mathrm{res}}\) is finite and nonempty.

Now let
\[
x\in\mathcal F_{m+1}.
\]
By Theorem~\ref{thm:exit-layer-first-contraction}, \(x\) realizes some
first-contraction word
\[
\eta=\omega d
\]
of length \(m+1\), with residual parent
\[
\omega\in\mathcal W_m^{\mathrm{res}}.
\]
Therefore
\[
c(\eta)
=
3c(\omega)+2^{D(\omega)}
<
K_m.
\]

Since \(\eta\) is a first-contraction word,
\[
D(\eta)\ge B(m+1),
\]
and hence
\[
3^{m+1}<2^{D(\eta)}.
\]
Thus
\[
2^{D(\eta)}-3^{m+1}
\]
is a positive integer and therefore is at least \(1\).

If \(x\ge K_m\), then
\[
c(\eta)<x,
\]
and consequently
\[
\frac{c(\eta)}
     {2^{D(\eta)}-3^{m+1}}
\le c(\eta)<x.
\]
By Proposition~\ref{prop:affine-descent-threshold},
\[
F_\eta(x)<x.
\]
Because \(x\) realizes \(\eta\),
\[
F_\eta(x)=f^{m+1}(x),
\]
so
\[
f^{m+1}(x)<x.
\]
\end{proof}

\begin{corollary}[Finite Non-Descent on Each Exit Layer]
\label{cor:finite-layer-nondescent}

For every \(m\ge1\), the set
\[
\mathcal E_{m+1}
:=
\left\{
x\in\mathcal F_{m+1}:
f^n(x)\ge x
\text{ for every }n\ge0
\right\}
\]
is finite.
\end{corollary}

\begin{proof}
By Theorem~\ref{thm:uniform-layer-descent}, every
\[
x\in\mathcal F_{m+1}
\]
with \(x\ge K_m\) satisfies
\[
f^{m+1}(x)<x.
\]
Hence every member of \(\mathcal E_{m+1}\) lies below \(K_m\).
\end{proof}

\begin{remark}
This uniformity is essential. A fixed first-contraction layer may contain
infinitely many first-contraction words, because the total division count
at the terminal depth may exceed the corresponding dyadic threshold by an
arbitrarily large amount. The finiteness of the exceptional set therefore
does not follow by taking an infinite union of the finite exceptional sets
associated with individual cylinders.

Instead, all children at a fixed layer inherit their affine constants from
the finitely many residual parents. This is what produces a single cutoff
\(K_m\) for the entire layer.
\end{remark}

% ------------------------------------------------------------
\subsection{Density Zero of Residual Non-Descent}
% ------------------------------------------------------------

Define the initially residual non-descent set
\[
\mathcal N_{\mathrm{res}}
:=
\left\{
x\in\mathcal R_1:
f^n(x)\ge x
\text{ for every }n\ge0
\right\}.
\]

The finite-stage forest decomposition gives, for every \(k\ge1\),
\[
\mathcal R_1
=
\left(
\bigsqcup_{j=2}^{k+1}\mathcal F_j
\right)
\sqcup
\mathcal R_{k+1}.
\]
Consequently, an initially residual integer that never descends either
survives to residual depth \(k+1\), or belongs to one of the finitely many
non-descending exceptional sets from the preceding exit layers.

\begin{proposition}[Finite-Depth Non-Descent Cover]
\label{prop:finite-depth-nondescent-cover}

For every \(k\ge1\), there exists a finite set
\[
E_k\subseteq\mathbb N
\]
such that
\[
\boxed{
\mathcal N_{\mathrm{res}}
\subseteq
\mathcal R_{k+1}\cup E_k.
}
\]
\end{proposition}

\begin{proof}
Set
\[
E_k
:=
\bigcup_{j=2}^{k+1}\mathcal E_j.
\]
Each \(\mathcal E_j\) is finite by
Corollary~\ref{cor:finite-layer-nondescent}, and the union contains only
finitely many layers, so \(E_k\) is finite.

Now let
\[
x\in\mathcal N_{\mathrm{res}}.
\]
If
\[
x\in\mathcal R_{k+1},
\]
there is nothing to prove. Otherwise, the finite-stage forest decomposition
places \(x\) in some
\[
\mathcal F_j,
\qquad
2\le j\le k+1.
\]
Since \(x\) never descends, it belongs to
\[
\mathcal E_j\subseteq E_k.
\]
\end{proof}

\begin{theorem}[Initially Residual Non-Descent Has Density Zero]
\label{thm:residual-nondescent-density-zero}

The set \(\mathcal N_{\mathrm{res}}\) has natural density zero:
\[
\boxed{
\operatorname{dens}(\mathcal N_{\mathrm{res}})=0.
}
\]
\end{theorem}

\begin{proof}
Let
\[
A(N)
:=
\#\bigl(\mathcal N_{\mathrm{res}}\cap[0,N)\bigr).
\]
Fix \(\varepsilon>0\).

By residual mass collapse, choose \(k\) such that
\[
M_{k+1}<\frac{\varepsilon}{2}.
\]
By Proposition~\ref{prop:finite-depth-nondescent-cover},
\[
\mathcal N_{\mathrm{res}}
\subseteq
\mathcal R_{k+1}\cup E_k,
\]
where \(E_k\) is finite. Writing
\[
C_k:=\#E_k,
\]
we obtain for every \(N\)
\[
A(N)
\le
\#\bigl(\mathcal R_{k+1}\cap[0,N)\bigr)+C_k.
\]
Hence
\[
\frac{A(N)}{N}
\le
\frac{\#(\mathcal R_{k+1}\cap[0,N))}{N}
+
\frac{C_k}{N}.
\]

The depth \(k+1\) is now fixed. Since
\[
\operatorname{dens}(\mathcal R_{k+1})
=
M_{k+1},
\]
the first term tends to \(M_{k+1}<\varepsilon/2\), while
\[
\frac{C_k}{N}\longrightarrow0.
\]
Therefore, for all sufficiently large \(N\),
\[
\frac{A(N)}{N}<\varepsilon.
\]
Since \(\varepsilon>0\) was arbitrary,
\[
\frac{A(N)}{N}\longrightarrow0.
\]
\end{proof}

\begin{remark}
The order of the limits is important. A finite residual depth is chosen
first, using
\[
M_m\longrightarrow0.
\]
Only after that depth has been fixed is the interval cutoff
\[
N\longrightarrow\infty
\]
taken.

Thus the argument uses neither countable additivity of natural density nor
an interchange of symbolic-depth and interval-size limits.
\end{remark}

% ------------------------------------------------------------
\subsection{The Initial Residual Congruence Class}
% ------------------------------------------------------------

The abstract initial residual population has a simple arithmetic
description.

\begin{theorem}[Initial Residual Class]
\label{thm:initial-residual-mod-four}

For every \(x\in\mathbb N\),
\[
\boxed{
x\in\mathcal R_1
\quad\Longleftrightarrow\quad
x\equiv3\pmod4.
}
\]
\end{theorem}

\begin{proof}
The unique residual word of length one is
\[
(1).
\]
Hence
\[
x\in\mathcal R_1
\]
if and only if \(x\) realizes the one-symbol word \((1)\), equivalently,
\[
v_2(3x+1)=1.
\]
By the exact valuation criterion,
\[
3x+1\equiv2\pmod4.
\]
This is equivalent to
\[
x\equiv3\pmod4.
\]
\end{proof}

Thus the residual forest begins with exactly one quarter of the natural
numbers. The remaining congruence classes are elementary.

\begin{proposition}[Immediate Descent Outside the Initial Residual Class]
\label{prop:immediate-descent-outside-residual}

The following hold.

\begin{enumerate}
\item Every positive even integer \(x\) satisfies
\[
T(x)<x,
\]
where \(T\) denotes the ordinary Collatz map.

\item Every \(x>1\) satisfying
\[
x\equiv1\pmod4
\]
satisfies
\[
f(x)<x.
\]
\end{enumerate}
\end{proposition}

\begin{proof}
If \(x\) is positive and even, then
\[
T(x)=\frac{x}{2}<x.
\]

Now suppose
\[
x>1,
\qquad
x\equiv1\pmod4.
\]
Then
\[
4\mid 3x+1,
\]
so
\[
v_2(3x+1)\ge2.
\]
Consequently,
\[
f(x)
=
\frac{3x+1}{2^{v_2(3x+1)}}
\le
\frac{3x+1}{4}.
\]
Since \(x>1\),
\[
3x+1<4x,
\]
and therefore
\[
f(x)<x.
\]
\end{proof}

\begin{remark}
Among positive integers, the only value not covered by the preceding
immediate-descent alternatives is \(1\), together with the congruence class
\[
x\equiv3\pmod4,
\]
which is precisely the initial residual class \(\mathcal R_1\).

In this sense, the residual-forest argument resolves the single nontrivial
quarter of the positive integers left after the elementary congruence
reduction.
\end{remark}

% ------------------------------------------------------------
\subsection{Transfer to the Ordinary Collatz Map}
% ------------------------------------------------------------

Define the ordinary Collatz map
\[
T:\mathbb N_{>0}\to\mathbb N_{>0}
\]
by
\[
T(x)
=
\begin{cases}
x/2, & x\equiv0\pmod2,\\[4pt]
3x+1, & x\equiv1\pmod2.
\end{cases}
\]

For an odd starting value \(x\), one fully accelerated step consists of one
ordinary odd Collatz step followed by exactly
\[
d(x)=v_2(3x+1)
\]
successive divisions by \(2\). Hence
\[
f(x)
=
T^{\,d(x)+1}(x).
\]

Iterating this observation gives the following bridge.

\begin{theorem}[Accelerated Iterates Occur Along the Ordinary Orbit]
\label{thm:accelerated-to-ordinary-iterate}

Let \(x\in\mathbb N_{>0}\) be odd. For every \(n\ge0\), there exists
\(m\ge0\) such that
\[
\boxed{
T^m(x)=f^n(x).
}
\]
\end{theorem}

\begin{proof}
Proceed by induction on \(n\).

For \(n=0\), take \(m=0\).

Suppose the result holds after \(n\) accelerated steps. One additional
fully accelerated step from the odd value \(f^n(x)\) is obtained by one
ordinary odd step followed by exactly
\[
d\!\left(f^n(x)\right)
\]
ordinary halving steps. Hence the endpoint \(f^{n+1}(x)\) is reached after
finitely many additional ordinary Collatz steps. Concatenating the two
finite ordinary segments proves the result.
\end{proof}

\begin{corollary}[Accelerated Descent Implies Ordinary Descent]
\label{cor:accelerated-implies-ordinary-descent}

Let \(x\in\mathbb N_{>0}\) be odd. If
\[
f^n(x)<x
\]
for some \(n\), then
\[
T^m(x)<x
\]
for some \(m\).
\end{corollary}

\begin{proof}
By Theorem~\ref{thm:accelerated-to-ordinary-iterate}, choose \(m\) such that
\[
T^m(x)=f^n(x).
\]
Then
\[
T^m(x)<x.
\]
\end{proof}

% ------------------------------------------------------------
\subsection{Density-One Descent}
% ------------------------------------------------------------

We can now assemble the argument.

Define the ordinary eventual-descent set
\[
\mathcal D
:=
\left\{
x\in\mathbb N_{>0}:
\exists m\ge0,\;
T^m(x)<x
\right\},
\]
and let
\[
\mathcal N
:=
\mathbb N_{>0}\setminus\mathcal D
\]
denote the ordinary non-descent set.

\begin{theorem}[Density-One Eventual Descent]
\label{thm:density-one-eventual-descent}

The set of positive integers admitting eventual descent under the ordinary
Collatz map has natural density one:
\[
\boxed{
\operatorname{dens}(\mathcal D)=1.
}
\]

Equivalently,
\[
\boxed{
\lim_{N\to\infty}
\frac{
\#\left\{
x\in\mathbb N_{>0}:
x<N,\;
\exists m\ge0,\;
T^m(x)<x
\right\}
}{N}
=
1.
}
\]
\end{theorem}

\begin{proof}
We first show that the ordinary non-descent set has density zero.

Let
\[
x\in\mathcal N
\]
with
\[
x\ne1.
\]
Then \(x>1\).

The integer \(x\) cannot be even, since every positive even integer descends
after one ordinary Collatz step. Hence \(x\) is odd.

It also cannot satisfy
\[
x\equiv1\pmod4,
\]
because Proposition~\ref{prop:immediate-descent-outside-residual} would give
\[
f(x)<x,
\]
and Corollary~\ref{cor:accelerated-implies-ordinary-descent} would then give
ordinary Collatz descent.

Therefore
\[
x\equiv3\pmod4.
\]
By Theorem~\ref{thm:initial-residual-mod-four},
\[
x\in\mathcal R_1.
\]

Moreover, \(x\) cannot descend under any accelerated iterate. If
\[
f^n(x)<x
\]
for some \(n\), then accelerated descent would again imply ordinary descent,
contrary to \(x\in\mathcal N\).

Hence
\[
x\in\mathcal N_{\mathrm{res}}.
\]
We have therefore proved
\[
\boxed{
\mathcal N
\subseteq
\mathcal N_{\mathrm{res}}\cup\{1\}.
}
\]

By Theorem~\ref{thm:residual-nondescent-density-zero},
\[
\operatorname{dens}(\mathcal N_{\mathrm{res}})=0.
\]
Moreover, for every \(N\ge1\), the inclusion above gives
\[
\#\bigl(\mathcal N\cap[0,N)\bigr)
\le
\#\bigl(\mathcal N_{\mathrm{res}}\cap[0,N)\bigr)+1.
\]
Therefore
\[
0
\le
\frac{\#\bigl(\mathcal N\cap[0,N)\bigr)}{N}
\le
\frac{\#\bigl(\mathcal N_{\mathrm{res}}\cap[0,N)\bigr)}{N}
+\frac1N.
\]
Both terms on the right tend to \(0\), so by the squeeze theorem,
\[
\operatorname{dens}(\mathcal N)=0.
\]

For every \(N\ge1\), descent and non-descent partition the positive integers
below \(N\):
\[
\{1,2,\ldots,N-1\}.
\]
Consequently,
\[
\#(\mathcal D\cap[0,N))
+
\#(\mathcal N\cap[0,N))
=
N-1.
\]
Dividing by \(N\) gives
\[
\frac{\#(\mathcal D\cap[0,N))}{N}
=
1-\frac1N
-
\frac{\#(\mathcal N\cap[0,N))}{N}.
\]
The last two terms tend to \(0\). Hence
\[
\operatorname{dens}(\mathcal D)=1.
\]
\end{proof}

\begin{remark}[Scope of the Theorem]
The theorem asserts that for a set of positive integers of natural density
one, some ordinary Collatz iterate is strictly smaller than the starting
value:
\[
\exists m,\qquad T^m(x)<x.
\]

It does not assert that every positive starting value has this property,
nor does it assert that every Collatz trajectory reaches \(1\). A
density-zero exceptional set may still be nonempty or infinite.
\end{remark}


% ============================================================
\section{Conclusion}
% ============================================================

This paper develops a symbolic--arithmetic framework for finite trajectories
of the fully accelerated Collatz map and uses it to prove density-one
eventual descent for the ordinary Collatz map. Division words encode exact
sequences of \(2\)-adic valuations, while the affine representation
\[
F_\omega(x)
=
\frac{3^{|\omega|}x+c(\omega)}{2^{D(\omega)}}
\]
translates that symbolic data into explicit arithmetic constraints.

Exact realizability requires one additional dyadic parity condition beyond
formal integrality. For every admissible division word \(\omega\), its
realizations form a unique residue class modulo
\[
2^{D(\omega)+1}.
\]
These realization classes therefore form dyadic cylinders of natural density
\[
2^{-(D(\omega)+1)},
\]
with symbolic prefix extension corresponding exactly to cylinder
containment.

The affine slope separates symbolic prefixes according to whether it remains
above \(1\) or has become subunit. This leads to the residual forest:
residual words are those whose nonempty prefixes remain above unit affine
slope, while first-contraction words cross below unit slope for the first
time at their terminal position. At each finite depth the residual family is
finite, and the corresponding realization sets satisfy the exact
decomposition
\[
\mathcal R_m
=
\mathcal F_{m+1}
\sqcup
\mathcal R_{m+1}.
\]
Thus the unresolved population is partitioned without loss or duplication
between first-contraction exits and continued residual survival.

The quantitative core of the argument is the decay of the residual mass.
An auxiliary weight
\[
\left(\frac38\right)^{D(\omega)}
\]
gives a uniform one-step contraction bound \(3/5\). Comparison with the
true dyadic cylinder mass, together with the integer threshold estimate
\[
B(m+5)\le B(m)+8,
\]
gives a strict contraction of the residual-mass envelope over five steps.
Consequently,
\[
M_m=\operatorname{dens}(\mathcal R_m)\longrightarrow0.
\]

Affine contraction alone does not force every realization of a fixed word
to descend, because the affine constant may obstruct descent for finitely
many small starting values. The residual forest resolves this uniformly:
at each first-contraction depth, the finitely many residual parents yield a
single finite cutoff above which every realization in that exit layer
descends. Combining this finite-layer control with residual mass collapse
shows that the set of initially residual trajectories that never descend has
natural density zero.

The remaining congruence classes require no residual analysis. Positive even
integers descend immediately under the ordinary Collatz map, while odd
integers congruent to \(1\pmod4\), apart from the trivial value \(1\),
descend after one fully accelerated step. The only nontrivial initial class
is
\[
x\equiv3\pmod4,
\]
which is precisely the realization set \(\mathcal R_1\) of the root word
\((1)\) of the residual forest. Accelerated descent for an odd integer
occurs at an actual iterate of the ordinary Collatz map, so the accelerated
result transfers directly to ordinary eventual descent.

Hence
\[
\boxed{
\lim_{N\to\infty}
\frac{
\#\left\{
x\in\mathbb N_{>0}:
x<N,\;
\exists m\ge0,\;
T^m(x)<x
\right\}
}{N}
=
1.
}
\]

The proof developed here is organized through a sequence of exact
mathematical structures:
\[
\boxed{
\begin{aligned}
\text{division words}
&\longrightarrow
\text{affine realization}
\longrightarrow
\text{dyadic cylinders}
\\
&\longrightarrow
\text{residual forest}
\longrightarrow
\text{residual mass collapse}
\longrightarrow
\text{eventual descent}.
\end{aligned}
}
\]
Several ingredients underlying this framework---including affine
representations of finite Collatz trajectories, finite residue coding, and
density-one finite stopping---have classical antecedents in the work of
Terras and Everett. The formulation used here reorganizes these ideas in
terms of exact accelerated division counts and a residual-cylinder
decomposition, leading to the residual-mass argument above. At each stage
the symbolic description is tied to exact arithmetic rather than to a
heuristic trajectory model, and the final density argument uses only
finite-stage decompositions and ordinary natural-density limits.

The result does not prove that every Collatz trajectory descends below its
starting value, nor does it prove convergence of every trajectory to \(1\).
A density-zero exceptional set may remain nonempty or infinite. What is
proved is that eventual descent below the initial value holds for a set of
positive integers of natural density one.

\appendix

% ============================================================
\newpage
\section{Explicit Computation of the Affine Map}
\label{app:canonical-c}
% ============================================================

Each division word
\[
\omega=(d_0,\dots,d_{m-1})
\]
determines the formal affine map
\[
F_\omega(x)
=
\frac{3^m x+c(\omega)}{2^{D(\omega)}},
\qquad
D(\omega)=\sum_{i=0}^{m-1}d_i,
\]
where the affine constant \(c(\omega)\) is determined uniquely by symbolic
composition.

By Theorem~\ref{thm:explicit-affine-constant}, the constant is given
explicitly by
\[
c(\omega)
=
\sum_{j=0}^{m-1}
3^{m-1-j}
2^{\sum_{i=0}^{j-1}d_i}.
\]
This is an integer identity, not a choice of residue representative modulo
\(2^{D(\omega)}\). In particular, \(c(\omega)\) need not lie in the interval
\[
0\le c(\omega)<2^{D(\omega)}.
\]

For example, if
\[
\omega=(1,1),
\]
then
\[
c(\omega)=3+2=5,
\]
and therefore
\[
F_\omega(x)=\frac{9x+5}{4}.
\]
Although
\[
5\equiv1\pmod4,
\]
replacing \(5\) by \(1\) would change the affine map:
\[
\frac{9x+1}{4}\ne\frac{9x+5}{4}.
\]
Thus the affine constant must be retained as the actual integer produced by
symbolic composition rather than reduced modulo the denominator.

% ------------------------------------------------------------
\subsection{A Four-Step Example}
% ------------------------------------------------------------

Consider
\[
\omega=(1,1,1,4).
\]
Then
\[
m=4,
\qquad
D(\omega)=7,
\]
and
\[
\begin{aligned}
c(\omega)
&=
3^3
+
3^2 2^1
+
3\,2^{1+1}
+
2^{1+1+1}
\\
&=
27+18+12+8
\\
&=
65.
\end{aligned}
\]
Hence
\[
\boxed{
F_\omega(x)=\frac{81x+65}{128}.
}
\]

The formal-integrality condition is
\[
81x+65\equiv0\pmod{128},
\]
which gives
\[
x\equiv15\pmod{128}.
\]
Thus every integer of the form
\[
x=15+128k
\]
makes the formal affine expression integral, and
\[
F_\omega(15+128k)
=
81k+10.
\]

Formal integrality, however, is weaker than exact realization. By
Theorem~\ref{thm:exact-realization-criterion}, exact realization is
equivalent to
\[
81x+65
\equiv
128
\pmod{256},
\]
which says precisely that the quotient remaining after removal of the
prescribed total power \(2^7\) is odd. Solving gives
\[
\boxed{
x\equiv143\pmod{256}.
}
\]

Thus the formal-integrality class
\[
15\pmod{128}
\]
splits into two classes modulo \(256\). Only
\[
143\pmod{256}
\]
realizes the division word \((1,1,1,4)\).

This distinction can also be seen directly from
\[
F_\omega(15+128k)=81k+10.
\]
The endpoint is odd exactly when \(k\) is odd. Writing
\[
k=2j+1
\]
gives
\[
x
=
15+128(2j+1)
=
143+256j,
\]
which is precisely the exact realization class.

For the canonical representative \(x=143\),
\[
F_\omega(143)
=
\frac{81\cdot143+65}{128}
=
91.
\]
The accelerated trajectory is
\[
143\longmapsto215\longmapsto323\longmapsto485\longmapsto91,
\]
with successive division counts
\[
(1,1,1,4).
\]

% ------------------------------------------------------------
\subsection{Direct Derivation by Symbolic Substitution}
% ------------------------------------------------------------

For completeness, consider a general four-entry division word
\[
\omega=(a,b,c,d).
\]
The formal symbolic relations are
\[
x_1=\frac{3x_0+1}{2^a},
\qquad
x_2=\frac{3x_1+1}{2^b},
\qquad
x_3=\frac{3x_2+1}{2^c},
\qquad
x_4=\frac{3x_3+1}{2^d}.
\]

Eliminating the intermediate variables gives
\[
2^{a+b+c+d}x_4
=
3^4x_0
+
3^3
+
3^2 2^a
+
3\,2^{a+b}
+
2^{a+b+c}.
\]
Therefore
\[
F_\omega(x)
=
\frac{
3^4x+
3^3+
3^2 2^a+
3\,2^{a+b}+
2^{a+b+c}
}{
2^{a+b+c+d}
},
\]
and hence
\[
\boxed{
c(\omega)
=
3^3+
3^2 2^a+
3\,2^{a+b}+
2^{a+b+c}.
}
\]

This calculation illustrates how each successive ``\(+1\)'' term is
propagated through the later affine steps. Earlier contributions acquire
additional factors of \(3\), while the powers of \(2\) record the prescribed
divisions occurring before each contribution reaches the final numerator.

The example also illustrates the distinction used throughout the paper:
the affine constant is determined by symbolic composition, formal
integrality is a congruence modulo \(2^{D(\omega)}\), and, for admissible
words, exact realization is characterized by the strengthened congruence
modulo \(2^{D(\omega)+1}\).


% ============================================================
\clearpage
\bibliographystyle{unsrtnat}
\bibliography{references}

\end{document}